\documentclass[a4paper,11pt]{article}
\pdfoutput=1
\pdfminorversion=7
\usepackage[utf8]{inputenc}
\usepackage{CJKutf8}
\AtBeginDocument{\begin{CJK*}{UTF8}{gbsn}}
\AtEndDocument{\clearpage\end{CJK*}}
\usepackage{amsmath,amssymb,latexsym}
\usepackage{bm}
\usepackage{booktabs}
\usepackage[svgnames]{xcolor}
\usepackage{hyperref}
\hypersetup{colorlinks=true,linkcolor=blue,citecolor=blue,urlcolor=blue}
\setlength{\emergencystretch}{3em}

\title{dS IBP Package 技术笔记}
\author{基于 2401.00129 (dS IBP) 与 1703.10166 (SK Diagrammatics)}

\begin{document}
\sloppy
\maketitle
\tableofcontents
\newpage

% ============================================================
\section{实现门禁与 seed canonical 约定}
% ============================================================

本 package 的目标是生成任意拓扑、任意 massive/massless 混合情形的 dS IBP seed。实现顺序上，Kira 导出不是前端公式的一部分，必须等以下条件全部满足后才开放：
\begin{itemize}
  \item time-IBP 与 momentum-IBP 两类生成元都已经实现；
  \item building-block 导数项已经包含在 seed 中；
  \item 对每个 sector 先枚举离散 $n=0,1$ 状态，再作用生成元；
  \item 生成元作用后一旦出现 $n=2$ 或更高 Hankel 导数态，立刻用 EOM 递推消去；
  \item 输出 seed 中所有完整线的 $n$ 指标均已 canonical 到 $0/1$，massless 线保持双端点三槽指标包 $\{b_e,n_{e,1},n_{e,2}\}$ 并应用 quotient relations；
  \item 同一代表顶点对的 full lines 按共同 theta 生成 odd-subset contact；sector 只枚举这些事件可达的状态。
\end{itemize}

当前 020 以既有 loop 物理核心为只读基线，已实现 momentum seed 的传播子项、shrunk-line $bS$ 项、z/ISP 吸收、massive/massless building-block 导数与 EOM/quotient canonical 门禁；time seed 已包含顶点幂次、外部相位、双端点 regular 项、massive/massless theta-boundary shrink、目标 sector coincident 抵消与 sector prefactor normalization，并通过标准 \texttt{dSIBP`} context 提供原子与高层接口。020 只在 time-only 公开边界引入式~\eqref{directPureTimeJ}，不重写这些 producer。massive 导数与 shrink 分别只读取由 $A_T/W_T$ 编译的数据。Kira serializer 只转换经过这些门禁的后端中立线性系统，importer 只验证并读取外部完整 reduction；二者都不替代任何前端关系或运行 reduction。

% ============================================================
\section{h 函数及其性质}
% ============================================================

本节从 dS IBP 文献 (2401.00129) Eq.~17--21 提取 building block 函数 $h$ 的定义、微分方程和递推关系。这些公式是 EOM 约化和 IBP 生成的基础。

\subsection{质量参数 $\nu$ 的定义}

参数 $\nu$ 由 dS 标量场的质量 $m$ 和时空维度决定。对 $\text{dS}_{d+1}$ 中的标量场（边界维度为 $d$）：
\begin{equation}
  \nu^2 = \frac{m^2}{H^2} - \frac{d^2}{4},
  \label{nuDef}
\end{equation}
其中 $H$ 是 Hubble 常数。本 package 中 $d$ 缺省为 $3$（即 $\text{dS}_4$ 体时空），此时：
\begin{equation}
  \nu^2 = \frac{m^2}{H^2} - \frac{9}{4}.
  \label{nuDef3d}
\end{equation}

$\nu$ 的取值分为两类：
\begin{itemize}
  \item \textbf{纯实数 $\nu$}（重场，$m > \frac{d}{2}H$）：$\nu = \sqrt{m^2/H^2 - d^2/4} \in \mathbb{R}^+$。此时 $\nu^* = \nu$，h 函数的两类 building block $h^{(1)}$ 和 $h^{(2)}$ 使用相同阶数的 Hankel 函数。
  \item \textbf{纯虚数 $\nu$}（轻场，$m < \frac{d}{2}H$）：$\nu = i\mu$，其中 $\mu = \sqrt{d^2/4 - m^2/H^2} \in \mathbb{R}^+$。此时 $\nu^* = -\nu = -i\mu$，h 函数的两类 building block 使用不同阶数（$\nu$ 与 $\nu^*$）的 Hankel 函数。
\end{itemize}

\textbf{本 package 只考虑纯实数和纯虚数 $\nu$}，不处理一般复数 $\nu$ 的情况。对 $d=3$，分界质量为 $m = \frac{3}{2}H$。

注意：部分文献使用相反符号约定 $\nu^2 = d^2/4 - m^2/H^2$，此时轻场对应实数 $\nu$、重场对应虚数 $\nu$。本 package 采用式~\eqref{nuDef} 的约定。

\subsection{定义}

文献 Eq.~17 定义 building block 函数为 Hankel 函数乘以归一化因子：
\begin{equation}
  h^{(1,2)}(\nu, 0; -k\tau) \equiv (-k\tau)^{-\nu} H_\nu^{(1,2)}(-k\tau),
  \label{hdef0}
\end{equation}
其中 $H_\nu^{(1,2)}$ 是第一/二类 Hankel 函数，$\nu$ 是质量参数（对 dS scalar $\nu^2 = m^2/H^2 - d^2/4$），$k = |\boldsymbol{k}|$ 是动量模长，$\tau < 0$ 是共形时间。

一阶导数态定义为：
\begin{equation}
  h^{(1,2)}(\nu, 1; -k\tau) \equiv -\frac{1}{k} \partial_\tau h^{(1,2)}(\nu, 0; -k\tau).
  \label{hdef1}
\end{equation}
此定义使得 $n \in \{0, 1\}$ 标记导数阶数，$n=0$ 为基函数，$n=1$ 为一次导数。

\subsection{微分方程}

从 Hankel 方程导出，$h(\nu, 0; -k\tau)$ 满足（文献 Eq.~18）：
\begin{equation}
  h''(\nu, 0; -k\tau) + \frac{2\nu+1}{\tau} h'(\nu, 0; -k\tau) + k^2 h(\nu, 0; -k\tau) = 0,
  \label{hODE}
\end{equation}
其中撇号表示对 $\tau$ 的导数。

令 $\tilde{\tau} = -k\tau$，方程变为（文献 Eq.~19）：
\begin{equation}
  \partial_{\tilde{\tau}}^2 h(\nu, 0; \tilde{\tau}) + \frac{2\nu+1}{\tilde{\tau}} \partial_{\tilde{\tau}} h(\nu, 0; \tilde{\tau}) + h(\nu, 0; \tilde{\tau}) = 0.
  \label{hODEtilde}
\end{equation}

\subsection{递推关系}

时间导数作用于 $h$ 函数产生指标递推（文献 Eq.~20）：
\begin{align}
  \partial_\tau h(\nu, 0; -k\tau) &= -k\, h(\nu, 1; -k\tau), \label{dh0} \\
  \partial_\tau h(\nu, 1; -k\tau) &= -k \left[ \frac{2\nu+1}{k\tau} h(\nu, 1; -k\tau) - h(\nu, 0; -k\tau) \right]. \label{dh1}
\end{align}

式~\eqref{dh0} 是 $h(\nu,1)$ 的定义 \eqref{hdef1} 的直接结果。式~\eqref{dh1} 由 ODE \eqref{hODE} 导出：将 $h''$ 用 ODE 替换后整理。

动量导数的递推关系类似（文献 Eq.~21）：
\begin{align}
  \partial_k h(\nu, 0; -k\tau) &= -\tau\, h(\nu, 1; -k\tau), \label{dkh0} \\
  \partial_k h(\nu, 1; -k\tau) &= -\tau \left[ \frac{2\nu+1}{k\tau} h(\nu, 1; -k\tau) - h(\nu, 0; -k\tau) \right]. \label{dkh1}
\end{align}

\subsection{EOM 递推规则（指标语言）}

将式~\eqref{dh1} 翻译为指标操作。被积函数中 $h(\nu, n; -k\tau)$ 与 $(-\tau)^a q^{-b}$ 的乘积在 $\partial_\tau$ 下产生：

对 $n=1$ 态的 $\tau$ 导数：
\begin{equation}
  \partial_\tau [(-\tau)^a q^{-b} h(\nu, 1)] = (-\tau)^{a-1} q^{-b} \left[ -a\, h(\nu,1) + (2\nu+1) h(\nu,1) + k\tau\, h(\nu,0) \right].
\end{equation}

在 IBP 中令此式为零（全微分 = 边界项），得到 $n=2$ 的 EOM 递推：
\begin{equation}
  \boxed{h(\nu, 2; -k\tau) = \frac{2\nu+1}{k\tau} h(\nu, 1; -k\tau) - h(\nu, 0; -k\tau).}
  \label{EOMrecurrence}
\end{equation}

在指标语言中（$n_{e,a} = 2$ 的积分用 $n_{e,a} = 1, 0$ 的积分表示）：
\begin{align}
  J[\ldots, \{b_e, \ldots, 2, \ldots\}, \ldots] &= -(2\nu_e + 1) \cdot J[\ldots, \{b_e + 1, \ldots, 1, \ldots\}, \ldots]_{\text{with } a_v \to a_v \pm 1} \notag \\
  &\quad + 1 \cdot J[\ldots, \{b_e, \ldots, 0, \ldots\}, \ldots],
  \label{EOMindex}
\end{align}
其中 $a_v$ 的移位方向取决于端点：端点 1 ($v = u[e]$) 取 $a_v \to a_v - 1$，端点 2 ($v = v[e]$) 取 $a_v \to a_v + 1$。$b_e \to b_e + 1$（端点 1）或 $b_e \to b_e - 1$（端点 2）。

EOM 递推系数总结：
\begin{equation}
  c_1 = 2\nu + 1, \quad c_2 = 1.
  \label{EOMcoeffs}
\end{equation}

对裸 Hankel，若 $H_0=H_\nu(x)$、$H_1=\partial_xH_0$，Bessel 方程给
\begin{equation}
 H_2=-H_0-\frac{1}{x}H_1+\frac{\nu^2}{x^2}H_0.
 \label{bareHEOMrecurrence}
\end{equation}
因此它不能压成常数二元系数。007--010 的内置 H 分支以及 012 的 H preset 在指标语言中实现为
\begin{equation}
 J_H[n=2]=-J_H[n=0]-J_H[a_v-1,b_e+1,n=1]
 +\nu_e^2J_H[a_v-2,b_e+2,n=0].
 \label{bareHEOMindex}
\end{equation}

\subsection{$P,Q,T,W$ 输入与 $2\times2$ 编译}

012 输入下一 subsection 统一定义的 $P,Q$、基变换 $T$（缺省单位矩阵）和原始导数基底带完整归一化的 Wronskian $W$。初始化层由 $P,Q,T,W$ 编译 $A_T$ 与 $W_T$，再分别生成 IBP 使用的导数移位和 shrink 数据；IBP 层不重算这些原始输入。

编译后的 h/H 一阶矩阵应分别为
\begin{equation}
 A_h(x)=
 \begin{pmatrix}0&1\\-1&-(2\nu+1)/x\end{pmatrix},
 \qquad
 A_H(x)=
 \begin{pmatrix}0&1\\\nu^2/x^2-1&-1/x\end{pmatrix}.
 \label{hH2x2Presets}
\end{equation}
这里 $h_1=\partial_xh_0$、$H_1=\partial_xH_0$。第二个矩阵直接来自
$H_0''+x^{-1}H_0'+(1-\nu^2/x^2)H_0=0$。012 的缺省 h 直接取 $P_h=(2\nu+1)/x$、$Q_h=1$、$T=I_2$ 和完整 $W_h$，故 $W_T=W_h$；裸 H 使用独立 preset。两者及一般 $T$ 均由专项检查和独立 EOM expected 验证。

\subsection{Wronskian：$P,Q$ 标准型与 $T$ 变换}
\label{wronskianPQTransform}

统一取二阶方程
\begin{equation}
 f''+P(x)f'+Q(x)f=0,
 \qquad
 A_0=\begin{pmatrix}0&1\\-Q&-P\end{pmatrix},
 \label{PQStandardODE}
\end{equation}
其中 $P$ 是一阶导系数，$Q$ 是零阶系数。对两个独立解定义
$W=f_1f_2'-f_1'f_2$。直接代入式~\eqref{PQStandardODE}，
\begin{equation}
 W'=f_1f_2''-f_1''f_2=-P\bigl(f_1f_2'-f_1'f_2\bigr)=-PW,
 \qquad
 W=C\exp\!\left(-\int P\,dx\right).
 \label{AbelPQ}
\end{equation}
$Q$ 项在相减时抵消；因此 $P$ 决定 $W$ 的变量依赖，解的归一化决定常数 $C$。

令 $\mathbf F=(f,f')^T$，两个独立解组成基本解矩阵
$\Phi=(\mathbf F^{(1)},\mathbf F^{(2)})$，故 $W=\det\Phi$。若同一个可逆矩阵 $T(x)$ 作用于两个解，$\Phi_T=T\Phi$，则
\begin{align}
 W_T&=\det(T)W,
 & A_T&=T'T^{-1}+TA_0T^{-1}, \notag\\
 \frac{W_T'}{W_T}
 &=\frac{W'}{W}+\frac{d}{dx}\log\det T
 =\operatorname{tr}A_T.
 \label{WronskianTTransform}
\end{align}
这同时给出 Wronskian 的代数变换和微分检验。

裸 Hankel 方程有 $P_H=1/x$、$Q_H=1-\nu^2/x^2$，故 $W_H=C_H/x$。对
\begin{equation}
 \mathbf h=x^{-\nu}
 \begin{pmatrix}1&0\\-\nu/x&1\end{pmatrix}\mathbf H,
 \qquad \det T=x^{-2\nu},
\end{equation}
式~\eqref{WronskianTTransform} 立即给出 $W_h=x^{-2\nu}W_H\propto x^{-2\nu-1}$。

% ============================================================
\section{线缩并（Shrink）机制}
% ============================================================

本节从 dS IBP 文献 Eq.~24--28 推导时间 IBP 中 Heaviside 函数导数产生的 delta 缩并项。

\subsection{物理机制}

时间 IBP 算子 $d/d\tau_v$ 作用于被积函数时，对 Heaviside 函数 $\theta(\tau_{u[e]} - \tau_{v[e]})$ 求导产生 delta 函数：
\begin{equation}
  \partial_{\tau_v} \theta(\tau_{u[e]} - \tau_{v[e]}) = \pm \delta(\tau_{u[e]} - \tau_{v[e]}),
  \label{thetaderiv}
\end{equation}
其中符号取决于 $v$ 是 $u[e]$ 还是 $v[e]$。

delta 函数将两个顶点 $u[e]$ 和 $v[e]$ 粘合为同一时间点，线 $e$ 的传播子被「缩并」。文献 Eq.~24 给出：
\begin{equation}
  \int \hat{V}_i(\ldots;\tau_i) \delta(\tau_i - \tau_j) \hat{V}_j(\ldots;\tau_j) d\tau_i d\tau_j = \int \hat{V}_i(\ldots;\tau_i) \hat{V}_j(\ldots;\tau_i) d\tau_i,
  \label{deltacollapse}
\end{equation}
即对 $\tau_j$ 的积分被消除，所有 $\tau_j$ 替换为 $\tau_i$。

\subsection{缩并线的 h 函数乘积}

当线 $e$ 被缩并（$\tau_{u[e]} \to \tau_{v[e]} \equiv \tau_v$），该线的两个端点处的 $h$ 函数乘积需要求值。文献 Eq.~25--28 分析了四种情况（$h^{(1,2)}$ 的组合），其中两种为零（反对称抵消），两种非零。

非零情况给出 Wronskian 型表达式（文献 Eq.~27）：
\begin{equation}
  F(-k_e\tau_v) = h^{(1)}(\nu_e, 1; -k_e\tau_v) h^{(2)}(\nu_e, 0; -k_e\tau_v) - h^{(2)}(\nu_e, 1; -k_e\tau_v) h^{(1)}(\nu_e, 0; -k_e\tau_v).
  \label{Fdef}
\end{equation}

采用标准方向 $W_z[f,g]=fg'-gf'$。利用 $W[H_\nu^{(1)}(z), H_\nu^{(2)}(z)] = -4i/(\pi z)$ 以及 $h$ 的定义 \eqref{hdef0}--\eqref{hdef1}，计算：
\begin{align}
  h^{(1)}(\nu, 1; z) &= -\frac{1}{k} \partial_\tau [(-k\tau)^{-\nu} H_\nu^{(1)}(-k\tau)] \notag \\
  &= z^{-\nu}\left[{H_\nu^{(1)}}'-\frac{\nu}{z}H_\nu^{(1)}\right].
\end{align}
因此式~\eqref{Fdef} 的方向是标准 Wronskian 的负号，
\begin{equation}
 F=-z^{-2\nu}W_z[H_\nu^{(1)},H_{\nu^*}^{(2)}].
 \label{FminusWronskian}
\end{equation}

代入 Wronskian 并化简（详细代数略），得到文献 Eq.~28：
\begin{equation}
  \boxed{F(-k_e\tau_v) = \mathcal{N}(\nu_e) \cdot (-k_e\tau_v)^{-(2\nu_e + 1)},}
  \label{Fresult}
\end{equation}
其中 $\mathcal{N}(\nu_e) = e^{\pi \text{Im}[\nu_e]} \cdot 4i/\pi$ 是纯系数（不依赖 $\tau$ 或 $k$）。

\subsection{缩并对指标的影响}

式~\eqref{Fresult} 的关键结构是 $(-k_e\tau_v)^{-(2\nu_e+1)}$。展开为：
\begin{equation}
  (-k_e\tau_v)^{-(2\nu_e+1)} = (-1)^{-(2\nu_e+1)} \cdot k_e^{-(2\nu_e+1)} \cdot \tau_v^{-(2\nu_e+1)}.
  \label{Fexpand}
\end{equation}

三个因子各自的影响：

\textbf{(a) $\tau_v^{-(2\nu_e+1)}$ 因子}：被吸收进合并顶点的时间幂次。原被积函数含 $(-\tau_v)^{a_u+a0_u+a_v+a0_v}$，乘以 $\tau_v^{-(2\nu_e+1)}$ 后，物理幂次移位为 $-(2\nu_e+1)$。在 package 的整数指标/零点分解中，整数部分 $-1$ 进入指标，非整数部分 $-2\nu_e$ 进入零点：
\begin{equation}
  \boxed{a_{\text{merged}} = a_u + a_v - 1,\qquad a0_{\text{merged}} = a0_u + a0_v - 2\nu_e.}
  \label{ashift}
\end{equation}

\textbf{(b) $k_e^{-(2\nu_e+1)}$ 因子}：这是动量模长的幂次。当线 $e$ 携带圈动量时，$k_e = |\boldsymbol{Q}_e|$ 依赖圈动量。在动量 IBP ($q_l^\mu \partial/\partial q_l^\mu$) 中：
\begin{equation}
  q_l^\mu \frac{\partial}{\partial q_l^\mu} k_e^{-(2\nu_e+1)} = -(2\nu_e+1) \frac{q_l \cdot Q_e}{k_e^2} k_e^{-(2\nu_e+1)}.
  \label{kshift}
\end{equation}
在 package 的整数指标/零点分解中，物理幂次 $-(2\nu_e+1)$ 对应缩并线分母指标的整数移位 $+1$，非整数部分 $+2\nu_e$ 进入 $bS0_e$：
\begin{equation}
  \boxed{bS_e = b_e + 1,\qquad bS0_e = b0_e + 2\nu_e.}
  \label{bSdef}
\end{equation}

\textbf{(c) $(-1)^{-(2\nu_e+1)}$ 因子}：纯相位，合并入总系数。

\subsection{缩并公式汇总}

时间 IBP 中线 $e$ 缩并的完整效果：

\begin{enumerate}
  \item 线 $e$ 从完整态 $\{b_e, n_{e,1}, n_{e,2}\}$ 变为缩并态 $\{bS_e\}$，整数指标为 $bS_e=b_e+1$；非整数部分进入 $bS0_e=b0_e+2\nu_e$。
  \item 合并顶点时间幂次的整数指标为 $a_{\text{merged}} = a_u + a_v - 1$；非整数部分进入 $a0_{\text{merged}}=a0_u+a0_v-2\nu_e$。
  \item 总系数获得因子 $\mathcal{N}(\nu_e)$；纯相位按初始化 convention 吸收到 prefactor 规则。
  \item 缩并线的 $h$ 函数消失，不再有 EOM 递推和 $n$ 指标。
  \item 在 sub-sector 的动量 IBP 中，当前缩并线 pack 的第一项 $bS_e$ 作为分母幂次指标贡献参与。
\end{enumerate}

条件：$n_{e,1} + n_{e,2} = 1$（一个端点 $n=0$，另一个 $n=1$）。$n_{e,1} + n_{e,2} = 0$ 或 $2$ 时缩并不产生（$F = 0$ 或对应不同机制）。

\subsection{Hankel 函数与 h 函数的缩并比较}

\subsubsection{h 函数的精确定义（文献 Eq.~17）}

文献 Eq.~17 的定义为（注意 $h^{(2)}$ 使用 $\nu^*$ 阶）：
\begin{align}
  h^{(1)}(\nu, 0; -k\tau) &\equiv (-k\tau)^{-\nu} H_\nu^{(1)}(-k\tau), \\
  h^{(2)}(\nu, 0; -k\tau) &\equiv (-k\tau)^{-\nu} H_{\nu^*}^{(2)}(-k\tau).
  \label{hdef_explicit}
\end{align}

\textbf{关键细节}：$h^{(2)}$ 的 Hankel 函数阶为 $\nu^*$（复共轭），不是 $\nu$。这源于 SK 形式论中 $(+)$ 顶点用 $H_\nu^{(1)}$、$(-)$ 顶点用 $H_{\nu^*}^{(2)}$ 的约定。$\nu$ 为实数时 $\nu^*=\nu$，两者等价；$\nu$ 为复数时此区别至关重要。

\subsubsection{Wronskian 的精确计算}

缩并因子为 $F = h^{(1)}(\nu,1) h^{(2)}(\nu,0) - h^{(2)}(\nu,1) h^{(1)}(\nu,0)$（文献 Eq.~27）。令 $z = -k\tau > 0$：

\begin{align}
  F &=-z^{-2\nu}\left[H_\nu^{(1)}{H_{\nu^*}^{(2)}}'-H_{\nu^*}^{(2)}{H_\nu^{(1)}}'\right]
     =-z^{-2\nu}W_z[H_\nu^{(1)},H_{\nu^*}^{(2)}].
\end{align}

这里出现的是 $W[H_\nu^{(1)}, H_{\nu^*}^{(2)}]$（\textbf{不同阶} $\nu$ 与 $\nu^*$），而非标准的 $W[H_\nu^{(1)}, H_\nu^{(2)}]$（同阶 $\nu$）。

对 $z > 0$ 实数，$H_{\nu^*}^{(2)}(z) = [H_\nu^{(1)}(z)]^*$，故：
\begin{equation}
  W[H_\nu^{(1)}, H_{\nu^*}^{(2)}] = H_\nu^{(1)} [H_\nu^{(1)}]'^* - [H_\nu^{(1)}] [H_\nu^{(1)}]'^* = -2i\,\text{Im}\!\left[H_\nu^{(1)}(z) {H_\nu^{(1)}}'(z)^*\right].
\end{equation}

此 Wronskian 的值依赖 $\nu$（包括 $\text{Im}[\nu]$），不是常数。

对比标准 Wronskian（同阶 $\nu$，与 $\nu$ 无关）：
\begin{equation}
  W[H_\nu^{(1)}, H_\nu^{(2)}] = -\frac{4i}{\pi z} \quad (\text{对任意 } \nu \text{ 成立}).
  \label{Wsameorder}
\end{equation}

文献 Eq.~28 给出 h 模式缩并因子的最终结果：
\begin{equation}
  F_h(-k\tau) = e^{\pi\text{Im}[\nu]} \frac{4i}{\pi} (-k\tau)^{-2\nu-1}.
  \label{Fresult_h}
\end{equation}

\subsubsection{H 模式（裸 Hankel）的缩并}

H 模式使用裸 Hankel 函数（无 $(-k\tau)^{-\nu}$ 归一化），但仍遵循 SK 约定 $H_\nu^{(1)}$/$H_{\nu^*}^{(2)}$：
\begin{equation}
  F_H(-k\tau) = -W_z[H_\nu^{(1)}, H_{\nu^*}^{(2)}].
\end{equation}

\textbf{数值验证结果}（见 \texttt{test/test\_wronskian.wl}，纯实数和纯虚数 $\nu$ 均测试，精度 $\sim 77$ 位）：
\begin{equation}
  \boxed{W[H_\nu^{(1)}, H_{\nu^*}^{(2)}] = -e^{\pi\text{Im}[\nu]} \frac{4i}{\pi z} \quad (\text{对纯实数和纯虚数 } \nu \text{ 成立}).}
  \label{WcrossOrderExact}
\end{equation}

因此 H 模式缩并因子为：
\begin{equation}
  F_H(-k\tau) = e^{\pi\text{Im}[\nu]} \frac{4i}{\pi} (-k\tau)^{-1}.
  \label{Fresult_H}
\end{equation}

$\tau$ 依赖为纯 $1/z$（整数幂次 $-1$），零点无移位。prefactor $= e^{\pi\text{Im}[\nu]} \frac{4i}{\pi}$，\textbf{与 h 模式相同}。

\subsubsection{$e^{\pi\text{Im}[\nu]}$ 因子的来源}

$e^{\pi\text{Im}[\nu]}$ 是以下两个因素\textbf{共同作用}的结果：

\begin{enumerate}
  \item \textbf{h 定义中 $h^{(2)}$ 使用 $H_{\nu^*}^{(2)}$}：使 Wronskian 变为 $W[H_\nu^{(1)}, H_{\nu^*}^{(2)}]$（$\nu$-依赖），而非 $W[H_\nu^{(1)}, H_\nu^{(2)}] = -4i/(\pi z)$（$\nu$-无关）。
  \item \textbf{归一化因子 $(-k\tau)^{-\nu}$}：两个 h 函数的归一化合并为 $z^{-2\nu}$，与 Wronskian 组合后产生最终的 $(-k\tau)^{-2\nu-1}$ 幂次。
\end{enumerate}

若 $h^{(2)}$ 改用 $H_\nu^{(2)}$（同阶），Wronskian 变为 $\nu$-无关的 $-4i/(\pi z)$，但 h 函数不再与 SK 顶点自然对应。

\subsubsection{缩并比较汇总}

\begin{itemize}
  \item \textbf{h 模式}：$F_h = e^{\pi\text{Im}[\nu]} \frac{4i}{\pi} (-k\tau)^{-2\nu-1}$。幂次 $-(2\nu+1)$ 分解为整数 $-1$（进入指标）和非整数 $-2\nu$（进入零点）。prefactor $= e^{\pi\text{Im}[\nu]} \frac{4i}{\pi}$。
  \item \textbf{H 模式}：$F_H = e^{\pi\text{Im}[\nu]} \frac{4i}{\pi} (-k\tau)^{-1}$（式~\eqref{Fresult_H}）。$\tau$ 依赖为纯 $1/z$（整数幂次 $-1$），零点无移位。prefactor $= e^{\pi\text{Im}[\nu]} \frac{4i}{\pi}$，\textbf{与 h 模式相同}。
  \item \textbf{无质量}：没有 Hankel Wronskian 型缩并；但双 theta 的 $n=1$ 端点导数仍产生 delta 缩并，第一/第二端点系数为 $-2/+2$，且 $bS_e=b_e$、prefactor $=1$。详见下一节。
\end{itemize}

\textbf{关键结论}：h 模式和 H 模式的缩并 prefactor 相同，均为 $e^{\pi\text{Im}[\nu]} \frac{4i}{\pi}$。差异仅在于 $\tau$ 幂次：h 模式为 $(-k\tau)^{-2\nu-1}$（含非整数部分 $-2\nu$，进入零点），H 模式为 $(-k\tau)^{-1}$（纯整数，零点无移位）。

% ============================================================
\section{Massless 双 theta 的 time-IBP 缩并与抵消}

本 package 对纯 massless 与 massive/massless 混合积分族都采用逐传播子的双 theta 合并路线。对一条 masslessFull 线，\texttt{lineData["endpoints"]=\{u,v\}} 是有序数据。令 $\Delta=\tau_u-\tau_v$，并对 $++/--$ 分别取 $\sigma=+1/-1$，先定义二维物理基
\begin{align}
 M_0 &=
 \theta(\Delta)e^{-i\sigma q\Delta}
 +\theta(-\Delta)e^{+i\sigma q\Delta},\\
 M_1 &=
 -\theta(\Delta)e^{-i\sigma q\Delta}
 +\theta(-\Delta)e^{+i\sigma q\Delta}.
 \label{masslessM01Tech}
\end{align}
交换端点时 $M_0$ 不变而 $M_1$ 变号。017 的公开 pack 保留双端点标签，并以
\begin{equation}
 F_{00}=M_0,\qquad F_{10}=M_1,\qquad
 F_{01}=-F_{10},\qquad F_{11}=-F_{00}
 \label{masslessDerivativeCanonicalTech}
\end{equation}
定义 quotient；动量幂变化另由 $b_e$ shift 保存。因此正式积分保存 $\{b_e,n_{e,1},n_{e,2}\}$，两个端点指标都只取 $0,1$，并以 $n_{e,2}\to0$ 为 canonical 方向；不建立 massless $n=2$ 状态。

式~\eqref{masslessDerivativeCanonicalTech} 是已抽出每次端点导数 $\pm i\sigma q$ 的无量纲 quotient，不是裸时间导数关系本身。原始 kernel 满足 $\{11\}=+q^2\{00\}$：cycle line 的两个 $q$ 由两次 $b_e\to b_e-1$ 吸收，fixed/non-loop line 则由导数算符显式乘两次结构化模长参数。因而 source seed 不需要同时生成四个双端点态；\texttt{DSSeeds} 只生成 $n_2=0$ 的 $\{00,10\}$ 两个代数代表，并在 metadata 中同时保存 raw 四态数、代表态数、删除数和方向 \texttt{n2ToZero}。导数若临时产生 $01/11$，立即按式~\eqref{masslessDerivativeCanonicalTech} 连同上述动量因子链约回。该缩减发生在离散 seed producer 层，不是连续指标 parity 筛点，也不得推迟到 \texttt{DSLinear} 去重。

\subsection{两个端点的 regular 与 theta-delta 项}

完整 kernel 的端点导数为
\begin{align}
 \partial_{\tau_u}M_n
 &=+i\sigma q\,M_{1-n}
   -2n\,\delta(\tau_u-\tau_v),\\
 \partial_{\tau_v}M_n
 &=-i\sigma q\,M_{1-n}
   +2n\,\delta(\tau_u-\tau_v).
 \label{masslessTimeRulesTech}
\end{align}
对一般公开态先应用式~\eqref{masslessDerivativeCanonicalTech}。在 canonical representative $\{b_e,n_e,0\}$ 上，regular 部分统一为
\begin{align}
 \tau_u:\quad
 \{b_e,n_e,0\}&\longrightarrow\{b_e-1,1-n_e,0\},
 &\text{系数 }+i\sigma,\\
 \tau_v:\quad
 \{b_e,n_e,0\}&\longrightarrow\{b_e-1,1-n_e,0\},
 &\text{系数 }-i\sigma.
\end{align}
连续在同一端点作用两次回到原 canonical representative，并给出 $-q^2$，即指标上为 $-\{b_e-2,n_e,0\}$。这一步直接由两次 $0\leftrightarrow1$ 翻转实现，不产生临时 $n_e=2$。

theta-delta 项在公开表示中只对奇数端点态 $n_{e,1}+n_{e,2}=1$ 存在。第一端点系数为 $-2$，第二端点为 $+2$，并触发
\begin{equation}
 \{b_e,n_{e,1},n_{e,2}\}\longrightarrow\{bS_e\},\qquad
 bS_e=b_e,\qquad bS0_e=b0_e.
 \label{masslessShrinkTech}
\end{equation}
 delta 合并两个时间变量；massless delta 没有额外 $1/(-\tau)$，故
\begin{equation}
 a_{\rm merged}=a_u+a_v,\qquad
 a0_{\rm merged}=a0_u+a0_v,
\end{equation}
且 massless 缩并没有 Hankel Wronskian prefactor，$\mathcal C_e=1$。这里 $bS_e=b_e$ 是强制区别：massive h/H 缩并才有 $bS_e=b_e+1$。

\subsection{同点抵消与目标 sector canonical}

若其它线的缩并使这条仍完整的 masslessFull 线的两个原端点映到同一个 active vertex，同一个 time 生成元必须同时作用两个端点。于是 regular 的 $+i\sigma/-i\sigma$ 项相消，theta-delta 的 $-2/+2$ 项相消，而且反对称态 $M_1(\tau,\tau)=0$。实现时不得用只返回第一个端点的搜索。

该规则也必须作用于 top-sector 方程中刚产生的 sub-sector 项。017 根据每个输出 $J$ 的 root-ordered full/shrunk pack pattern 重建目标 sector 的代表顶点映射，再判断 coincident 奇数端点态；若继续使用 source topology 的端点映射，会留下错误的 sub-sector 项。

masslessCross（$+-$ 或 $-+$）不含 theta，也没有活动离散态或 shrink。为保持固定三槽，公开 pack 写成 $\{b_e,0,0\}$ 或 fixed 的 $\{\texttt{"F"},0,0\}$；两个零不参加撒点。它的 time-IBP 只对两端点指数相位求导。

\subsection{共同-theta contact 到最终指标表示}

本节给出 package 采用的权威链路。两种分布实现的完整推导及等价条件见附录~\ref{app:commonThetaDerivation} 和附录~\ref{app:gaussianDerivation}。

设同一当前代表顶点对 $(u,v)$ 之间的 massive/massless full lines 构成 bundle $\mathcal B$，并统一用 $\Delta=\tau_u-\tau_v$ 表示时间差。逐线写成
\begin{equation}
 G_e=\theta(\Delta)A_e+\theta(-\Delta)B_e.
 \label{bundleSingleLine}
\end{equation}
共同 theta 先在整个乘积上合并：
\begin{equation}
 \prod_{e\in\mathcal B}G_e
 =\theta(\Delta)\prod_{e\in\mathcal B}A_e
 +\theta(-\Delta)\prod_{e\in\mathcal B}B_e.
 \label{bundleCommonThetaProduct}
\end{equation}
因此 time derivative 的奇异部分只有一个共同 delta，绝不产生 $\delta^k$：
\begin{equation}
 \left.\partial_\Delta\prod_{e\in\mathcal B}G_e\right|_{\rm contact}
 =\delta(\Delta)
 \left(\prod_{e\in\mathcal B}A_e-\prod_{e\in\mathcal B}B_e\right).
 \label{bundleUniqueBoundary}
\end{equation}

为保持逐线 pack，定义局部对称/跃变基底
\begin{equation}
 \mathsf J_e=\frac{A_e+B_e}{2},\qquad D_e=A_e-B_e.
 \label{bundleJDDefinition}
\end{equation}
这里 $\mathsf J_e$ 是一条线的等时对称因子，不是三槽积分 Head $J[\cdots]$。代数展开给出
\begin{equation}
 \prod_eA_e-\prod_eB_e
 =\sum_{\substack{\varnothing\ne S\subseteq\mathcal B\\ |S|\ {\rm odd}}}
 2^{1-|S|}\prod_{e\in S}D_e\prod_{e\notin S}\mathsf J_e.
 \label{commonThetaOddSubset}
\end{equation}

下一步必须把每个 $D_e$ 落到 line-local contact 数据，不能停在抽象 Wronskian。对 massiveFull 线，初始化先编译
\begin{equation}
 W_{T,e}(x_e)=\det T_e(x_e)\,W_e(x_e)
 =\sum_{\alpha}w_{e\alpha}x_e^{p_{e\alpha}},
 \qquad x_e=-q_e\tau,
 \label{WTLaurentContact}
\end{equation}
并把每个幂次唯一拆成
\begin{equation}
 p_{e\alpha}=-s_{e\alpha}-z_e,
 \qquad s_{e\alpha}\in\mathbb Z.
 \label{WTShiftSplit}
\end{equation}
同一条线所有 Laurent 项必须具有相同 $z_e$，否则函数系统不能编译。由于实际 shrink factor 是 $-W_{T,e}$，端点 $r$ 的 atomic contact 项为
\begin{equation}
 D_{e,r}\ \longmapsto\
 \delta_{n_{e,1}+n_{e,2},1}
 (-1)^{n_{e,r}+V_\pm}
 \sum_\alpha(-w_{e\alpha})x_e^{-s_{e\alpha}-z_e},
 \qquad V_+=1,\quad V_-=0.
 \label{massiveDToWT}
\end{equation}
这正是代码中的
\begin{center}
\texttt{compiledFunctionSystem["WT"]}
$\longrightarrow$
\texttt{shrinkTerms}
$\longrightarrow$
\texttt{thetaBoundaryAtomicTerms}。
\end{center}
对 h preset，$s_e=1,z_e=2\nu_e$ 且 $-W_{T,e}=\mathcal C_ex_e^{-2\nu_e-1}$；对 H preset，$s_e=1,z_e=0$ 且 $-W_{T,e}=\mathcal C_ex_e^{-1}$，其中
\begin{equation}
 \mathcal C_e=\frac{4i}{\pi}e^{\pi\operatorname{Im}\nu_e}.
\end{equation}
对 masslessFull，只有 $n_{e,1}+n_{e,2}=1$ 的 odd endpoint state 有 contact；第一/第二有序端点的系数为 $-2/+2$，并等价于 $s_e=z_e=0$ 的 atomic term，不调用 $W_T$。

现在选取式~\eqref{commonThetaOddSubset} 中一个奇数子集 $S$，并对每条 massive $e\in S$ 选定一个 Laurent 项 $\alpha_e$。记包含端点、SK 与 $-w_{e\alpha_e}$ 的完整 atomic 系数为 $c_{e\alpha_e}$。该项的总系数为
\begin{equation}
 C_{S,\boldsymbol\alpha}
 =2^{1-|S|}\prod_{e\in S}c_{e\alpha_e}.
 \label{bundleIndexCoefficient}
\end{equation}
若 bundle 两端在 sector 中对应 $a_u,a_v$，最终三槽积分的指标映射为
\begin{align}
 a_{uv}'&=a_u+a_v-\sum_{e\in S}s_{e\alpha_e},
 &
 a0_{uv}'&=a0_u+a0_v-\sum_{e\in S}z_e,
 \label{bundleMergedAShift}\\
 \{b_e,n_{e,1},n_{e,2}\}&\longrightarrow\{b_e+s_{e\alpha_e}\},
 &
 bS0_e&=b0_e+z_e,
 \qquad e\in S\text{ massive},
 \label{bundleMassivePackShift}\\
 \{b_e,n_{e,1},n_{e,2}\}&\longrightarrow\{b_e\},
 &
 bS0_e&=b0_e,
 \qquad e\in S\text{ massless}.
 \label{bundleMasslessPackShift}
\end{align}
因此这一项最终写成
\begin{equation}
 C_{S,\boldsymbol\alpha}\,
 J\!\left[
  \{\ldots,a_{uv}',\ldots\},
  \{\ldots,\{b_e+s_{e\alpha_e}\}_{e\in S},
          \text{full packs}_{e\notin S},\ldots\},
  \mathrm{ispList}
 \right].
 \label{bundleFinalJIndexForm}
\end{equation}
式~\eqref{bundleMergedAShift}--\eqref{bundleFinalJIndexForm} 中，所有选中线只触发一次顶点合并，但其整数与 zero-point shift 分别求和；各 shrunk line 的 $bS/bS0$ 仍逐线记录。未选中的 full lines 在合并后成为 coincident：massless odd endpoint state 立即置零，massiveFull 应用 $\{b,1,0\}\to\{b,0,1\}$；它们不再产生 theta boundary。sector 由这些非零 contact 事件做 BFS 得到，不能用 full-line 幂集代替。对某个离散端点态，若 atomic contact 的系数在结构上精确为零，必须在构造 Cartesian product 之前删除该 atomic choice；若所选线因此没有任何非零 choice，则该 odd subset 的贡献就是零。不得先构造 $0\,J$ 再送入 sector normalization，否则零项会触发不存在的 target-sector 归一化并把 \texttt{\$Failed} 污染到公开 \texttt{dtau}。

单传播子的对称等时值取 $\theta(0)=1/2$，对应式~\eqref{bundleJDDefinition} 的 $\mathsf J_e$。这只定义单线等时平均，不定义 $\delta\theta^m$；后者必须使用式~\eqref{bundleUniqueBoundary}，或附录~\ref{app:gaussianDerivation} 的统一正则化。
\section{维度参数约定}
% ============================================================

本 package 区分两个维度概念：

\begin{itemize}
  \item \textbf{PQ 维度 $d_{\text{PQ}}$}：出现在 $h$ 函数的 ODE 和 EOM 递推中（如式~\eqref{hODE} 的 $(2\nu+1)/\tau$ 项中的 $2\nu+1 = d-1$ 部分）。缺省值 $d_{\text{PQ}} = 3$。在 EOM 递推系数中已直接代入（$c_1 = 2\nu+1$ 不显式含 $d_{\text{PQ}}$）。
  \item \textbf{维正规化维度 $d$}：出现在动量 IBP 的维度项 $\text{dim} \cdot J$ 中。缺省值 $d = 3$（树图）或 $d = 3 - 2\epsilon$（圈图）。
\end{itemize}

两者独立，不做混用。

% ============================================================
\section{Building Block 参数体系}
% ============================================================

每条 massive 完整线 $e$ 在 012 中都带编译后的 \texttt{functionSystem}。推荐直接输入 $P,Q,T,W$；旧 $\text{bbType}_e$ 与二元 $\texttt{eomCoefficients}$ 只作为初始化兼容入口。若旧输入给出 $\{c_1,c_2\}$，初始化层解释为：

\begin{itemize}
  \item $P=c_1/x$、$Q=c_2$、$T=I_2$；
  \item $W=-\mathcal C x^{-c_1}$，整数 shrink shift 缺省取 $1$，zero-point shift 为 $c_1-1$。
\end{itemize}

内置缺省值为
\begin{align}
  \texttt{"h"} &\to \{P_h=(2\nu+1)/x,\;Q_h=1,\;T_h=I_2,\;W_h\}, \label{bbH} \\
  \texttt{"H"} &\to \{P_H=1/x,\;Q_H=1-\nu^2/x^2,\;T_H=I_2,\;W_H\}. \label{bbHankel}
\end{align}
其中 $A_H$ 是式~\eqref{hH2x2Presets}，由式~\eqref{bareHEOMindex} 编译为三项指标移位。缺省 h preset 取 $T=I_2$、$W=W_h$，所以 $W_T=W_h$；H preset 取 $T=I_2$、$W=W_H$。显式 $W_T$ 只用于核对 $W_T=\det(T)W$，不能覆盖编译结果。

% ============================================================
\section{指标零点体系}
% ============================================================

本 package 采用指标零点（zero-point）体系，将每个指标分解为整数部分（可变指标）和非整数/固定部分（零点）。这一分解使得 IBP 关系中的指标移位始终为整数，便于 Kira 等后端处理。

\subsection{基本定义}

对顶点 $v$ 的时间幂次指标 $a_v$（正幂次，即 $(-\tau_v)$ 的幂次）：
\begin{equation}
  \text{物理幂次} = a_v + a0_v, \quad \text{即被积函数含 } (-\tau_v)^{a_v + a0_v}.
  \label{adecomp}
\end{equation}

对线 $e$ 的动量幂次指标 $b_e$（分母幂次，即 $q_e$ 的幂次取负）：
\begin{equation}
  \text{物理幂次} = -(b_e + b0_e), \quad \text{即被积函数含 } q_e^{-(b_e + b0_e)}.
  \label{bdecomp}
\end{equation}

注意 $a$ 是正幂次（分子），$b$ 是分母幂次。物理幂次与指标的关系：$a$ 的物理幂次 $= +a$，$b$ 的物理幂次 $= -b$。

\subsection{零点缺省值}

零点缺省值由 building block 类型决定。验证依据：参考代码 \texttt{001 bubble\_ibp\_sym.m} 的 \texttt{reppara2N} 给出 $a0 \to 2\nu$, $b0 \to -2\nu$。

\begin{align}
  \texttt{"h" 模式:} \quad a0_v &= 2\nu_{\text{ref}}, & b0_e &= -2\nu_e, \label{zeroHmode} \\
  \texttt{"H" 模式:} \quad a0_v &= 0, & b0_e &= 0. \label{zeroHmodeH}
\end{align}

其中 $\nu_{\text{ref}}$ 是该顶点关联的线的质量参数（对 bubble：两线共享同一 $\nu$，故 $a0 = 2\nu$）。

\subsection{缩并时的零点移位}

当线 $e$ 缩并时，式~\eqref{Fresult} 的缩并因子 $(k_e \tau_v)^{-(2\nu_e+1)}$ 的幂次分解为整数部分和非整数部分：
\begin{equation}
  -(2\nu_e + 1) = \underbrace{-1}_{\text{整数}} + \underbrace{(-2\nu_e)}_{\text{非整数}}.
  \label{shrinkDecomp}
\end{equation}

\textbf{对 $a$ 零点}（$\tau$ 幂次）：合并顶点的零点为
\begin{equation}
  a0_{\text{merged}} = a0_u + a0_v + (-2\nu_e).
  \label{a0merged}
\end{equation}
整数部分 $-1$ 进入指标：$a_{\text{merged}} = a_u + a_v - 1$。

验证：参考代码 \texttt{repvar} 中 $a0R \to 2a0 - 2\nu = 2(2\nu) - 2\nu = 2\nu$。式~\eqref{a0merged} 给出 $a0_{\text{merged}} = 2\nu + 2\nu - 2\nu = 2\nu$。一致。

\textbf{对 $b$ 零点}（$q$ 幂次，分母幂次）：缩并线的零点为
\begin{equation}
  bS0_e = b0_e + 2\nu_e.
  \label{bS0def}
\end{equation}
注意符号：$b$ 的物理幂次 $= -b$，缩并因子的 $k$ 幂次为 $-(2\nu_e+1)$，对应物理 $q$ 幂次移位 $-(2\nu_e+1)$。由于物理幂次 $= -(b + b0)$，物理幂次移位 $-(2\nu_e+1)$ 对应 $b+b0$ 移位 $+(2\nu_e+1) = +1 + 2\nu_e$。整数部分 $+1$ 进入指标，非整数部分 $+2\nu_e$ 进入零点。

验证：参考代码 \texttt{repvar} 中 $b10R \to b0 + 2\nu = -2\nu + 2\nu = 0$。式~\eqref{bS0def} 给出 $bS0 = -2\nu + 2\nu = 0$。一致。

\textbf{剩余线}的 $b$ 零点不变：$b0_{\text{remaining}} = b0_e$。

\subsection{H 模式的缩并与零点}

H 模式使用裸 Hankel 函数 $H_\nu^{(1)}$/$H_{\nu^*}^{(2)}$（无 $(-k\tau)^{-\nu}$ 归一化），零点缺省为 0（式~\eqref{zeroHmodeH}）。

缩并反对称组合为 $F_H=-W_z[H_\nu^{(1)},H_{\nu^*}^{(2)}]$。对本 package 允许的纯实数和纯虚数 $\nu$，数值验证给出
\begin{equation}
  W[H_\nu^{(1)}, H_{\nu^*}^{(2)}] = -e^{\pi\text{Im}[\nu]} \frac{4i}{\pi z}.
\end{equation}
因此
\begin{equation}
  F_H = e^{\pi\text{Im}[\nu]} \frac{4i}{\pi}(-k\tau)^{-1}.
\end{equation}
整数移位部分与 h 模式相同（$a$-shift $= -1$, $b$-shift $= +1$），因为 Wronskian 的 $1/z$ 依赖给出 $(-k\tau)^{-1}$ 的整数幂次；H 模式无非整数 $-2\nu$ 幂次，因此零点无移位。

\subsection{缩并常数 prefactor}

每次缩并除产生指标移位和零点移位外，还产生一个常数 prefactor $\mathcal{C}_e$，乘入总体系数。此 prefactor 来自 Wronskian 计算中不依赖 $\tau$ 的部分。

\textbf{h 模式}（文献 Eq.~28）：
\begin{equation}
  \mathcal{C}_e^{(h)} = \frac{4i}{\pi} \cdot e^{\pi \text{Im}[\nu_e]}.
  \label{prefactorH}
\end{equation}

此 $e^{\pi\text{Im}[\nu]}$ 因子是以下两个因素\textbf{共同作用}的结果：
\begin{enumerate}
  \item h 定义中 $h^{(2)}$ 使用 $H_{\nu^*}^{(2)}$（而非 $H_\nu^{(2)}$），使 Wronskian 变为 $W[H_\nu^{(1)}, H_{\nu^*}^{(2)}] = -2i\,\text{Im}[H_\nu^{(1)} {H_\nu^{(1)}}'^*]$，此量依赖 $\nu$（包括 $\text{Im}[\nu]$）。
  \item 归一化因子 $(-k\tau)^{-\nu}$ 与 Wronskian 组合后，$\nu$-依赖部分整理为 $e^{\pi\text{Im}[\nu]}$。
\end{enumerate}

当 $\nu_e$ 为实数时 $\text{Im}[\nu_e] = 0$，prefactor 简化为 $4i/\pi$。当 $\nu_e = i\mu$ 为纯虚数（轻场）时，$e^{\pi\mu}$ 给出指数增强。

\textbf{H 模式}：缩并反对称组合为 $F_H=-W_z[H_\nu^{(1)},H_{\nu^*}^{(2)}]$。对纯实数和纯虚数 $\nu$，已数值验证
\begin{equation}
  W[H_\nu^{(1)}, H_{\nu^*}^{(2)}] = -e^{\pi\text{Im}[\nu]} \frac{4i}{\pi z}.
  \label{WcrossOrder}
\end{equation}
因此 H 模式的常数 prefactor 为
\begin{equation}
  \mathcal{C}_e^{(H)} = \frac{4i}{\pi} e^{\pi\text{Im}[\nu_e]},
\end{equation}
与 h 模式相同；差异只在于 H 模式的 $\tau$ 依赖为纯 $(-k\tau)^{-1}$，零点无移位。

注意：标准 Wronskian $W[H_\nu^{(1)}, H_\nu^{(2)}] = -4i/(\pi z)$（对任意 $\nu$ 成立）是\textbf{同阶} $\nu$ 的结果，与物理缩并中出现的 $W[H_\nu^{(1)}, H_{\nu^*}^{(2)}]$（\textbf{异阶} $\nu$ 与 $\nu^*$）不同。

\textbf{无质量传播子}：无质量标量在 dS 中的模式函数为 $e^{\pm ik\tau}$（指数函数，非 Hankel），其 Wronskian 为常数：
\begin{equation}
  W[e^{ik\tau}, e^{-ik\tau}] = -2ik.
\end{equation}
无 Hankel 缩并机制，故无 $\mathcal{C}_e$ prefactor（或等价地 $\mathcal{C}_e^{(\text{massless})} = 1$）。这里说的“无缩并机制”只指没有 Hankel Wronskian 型的 $h/H$ 塌缩 prefactor；massless 的 Heaviside boundary 仍按式~\eqref{commonThetaOddSubset} 与同顶点对的其它 full lines 一起处理，并保持逐线指标包 $\{b_e,n_{e,1},n_{e,2}\}$。

\textbf{多次缩并}：当多条线同时缩并时（如 2-line shrink），总 prefactor 为各线 prefactor 之积：
\begin{equation}
  \mathcal{C}_{\text{total}} = \prod_{e \in \text{shrunk}} \mathcal{C}_e.
  \label{prefactorTotal}
\end{equation}

在 package 实现中，$\mathcal{C}_e$ 作为每条缩并线的属性存储，在导出 Kira 输入时乘入 sub-sector 方程的系数。

% ============================================================
\section{多圈 IBP 生成：链式法则与复合算符}
% ============================================================

本节描述如何从拓扑输入自动生成圈动量 IBP 关系。核心思想是将 $\partial/\partial q_l^\mu$ 通过链式法则分解为对 $\xi_e$ 和 ISP 的导数，再与矢量 $v^\mu$ 缩并。

\subsection{链式法则分解}

每条内线 $e$ 的动量为 $Q_e^\mu = \sum_l c_{e,l}\, q_l^\mu + P_e^\mu$，其中 $c_{e,l}$ 为整数系数，$P_e^\mu$ 为外动量组合。定义 $\xi_e = |Q_e| = \sqrt{Q_e \cdot Q_e}$。

被积函数通过以下途径依赖圈动量 $q_l^\mu$：
\begin{enumerate}
  \item 传播子分母 $\xi_e^{-(b_e + b0_e)}$
  \item Building block $h(\nu_e, n_{e,a}; \xi_e)$
  \item ISP 分子因子 $\text{isp}_j^{n_j}$
  \item 指数相位 $e^{i P_v \cdot \tau_v}$（通常不含 $q_l$）
\end{enumerate}

链式法则给出：
\begin{equation}
  \frac{\partial}{\partial q_l^\mu} = \sum_e \frac{\partial \xi_e}{\partial q_l^\mu} \frac{\partial}{\partial \xi_e} + \sum_j \frac{\partial\, \text{isp}_j}{\partial q_l^\mu} \frac{\partial}{\partial\, \text{isp}_j}
\end{equation}

其中：
\begin{equation}
  \frac{\partial \xi_e}{\partial q_l^\mu} = c_{e,l} \frac{Q_{e,\mu}}{\xi_e}
\end{equation}

\subsection{复合 IBP 算符}

IBP 恒等式来自全微分在维正规化下积分为零：
\begin{equation}
  0 = \int \prod_l d^d q_l\; \frac{\partial}{\partial q_l^\mu} \left[ v^\mu \cdot F \right] = \int \prod_l d^d q_l\; \left[ (\partial \cdot v)\, F + v^\mu \frac{\partial F}{\partial q_l^\mu} \right]
\end{equation}

\textbf{复合 IBP 算符}定义为：
\begin{equation}
  \mathcal{O}_{l,v} = (\partial \cdot v) + v^\mu \frac{\partial}{\partial q_l^\mu}
\end{equation}

其中 $v^\mu$ 为可用动量的任意线性组合。代入链式法则：
\begin{equation}
  v^\mu \frac{\partial F}{\partial q_l^\mu} = \sum_e c_{e,l} \frac{v \cdot Q_e}{\xi_e} \frac{\partial F}{\partial \xi_e} + \sum_j (v \cdot \partial_{q_l} \text{isp}_j) \frac{\partial F}{\partial\, \text{isp}_j}
\end{equation}

每一项的作用：
\begin{itemize}
  \item $\frac{1}{\xi_e}\frac{\partial}{\partial \xi_e}$ 作用在传播子上时，先产生 $\xi_e^{-2}$，即在吸收 $v\cdot Q_e$ 前执行 $b_e\to b_e+2$；随后 $v\cdot Q_e$ 的 $z$/ISP 分解再决定各项的进一步指标移位。
  \item $\frac{\partial}{\partial \xi_e}$ 作用在 massive building block 上时，只读取该线最终 $A_T$ 编译成的 \texttt{derivativeTerms}；每个 $x=-\xi_e\tau$ Laurent 项同时给出目标基底态和整数 $a/b$ 移位，不再假设一般基底满足简单的 $n\to n+1$。
  \item $\frac{\partial}{\partial\, \text{isp}_j}$ 作用在 ISP 因子上：$n_{\text{isp}_j} \to n_{\text{isp}_j} - 1$
\end{itemize}

\subsection{完备 IBP 生成元集合}

参考 FIRE7 和 LiteRed 的方法，对 $L$ 圈积分，017 要求用户显式给出有序的 \texttt{loopExternalMomenta} 基 $\{k_1,\ldots,k_K\}$；程序不按名字或 line 首次出现自动选择角色。该基的平方 Gram 原子为 $x_{ij}=\operatorname{sp}(k_i,k_j)$，公开变量为 $\sqrt{x_{ij}}$。另一个显式列表 \texttt{independentExternalMomenta} 只声明实际无圈 line/phase 的独立模长；这些对象不计入 $K$、loop IBP 生成元或 ISP 空间，也不自动产生彼此交叉点积。

实际需要的 loop-external 空间由 affine routing 独立审计。若所有 line momentum 写成 $Q=Aq+r$，取一组满秩参考行 $A_R$ 并定义
\begin{equation}
 q'=A_Rq+r_R,\qquad r'=r-AA_R^{-1}r_R.
 \label{affineRoutingInvariant}
\end{equation}
$r'$ 的行空间及 ISP 中与圈动量耦合的外方向才是 shift-invariant 需求。例如 $\{q+k_1,q+k_1+k_2\}$ 只要求 $k_2$。程序把该需求与用户基比较以判断 exact/over/under，但不替用户改写输入。若声明过完备，原列表保留为 declared metadata，而实际用于 $K$、Gram 原子、$dqk$ 和 ISP 闭合的 $\texttt{effectiveLoopExternalMomenta}$ 取 $r'$ 需求空间的一组独立基。否则可吸收到 $q'$ 的共同平移方向会虚增 $q\cdot k$ 原子，并错误触发 ISP 不足。完备的 IBP 生成元为：

\begin{equation}
  \mathcal{O}_{l,v} = \frac{\partial}{\partial q_l^\mu} \cdot v^\mu, \quad l \in \{1,\ldots,L\}, \quad v^\mu \in \{q_1^\mu, \ldots, q_L^\mu, k_1^\mu, \ldots, k_K^\mu\}
\end{equation}

总数 $N_{\text{IBP}} = L(L + K)$，分解为：

\begin{itemize}
  \item \textbf{对角（scale）}：$v = q_l$，共 $L$ 个。$\partial \cdot q_l = d$。
  \item \textbf{交叉（cross）}：$v = q_m$（$m \neq l$），共 $L(L-1)$ 个。$\partial \cdot q_m = 0$（当 $m \neq l$）。
  \item \textbf{外动量（special）}：$v = k_j$，共 $LK$ 个。$\partial \cdot k_j = 0$。
\end{itemize}

\textbf{交叉 IBP 的必要性}：$L \geq 2$ 时，仅对角 IBP 不足以将所有积分约化到 master integrals。交叉 IBP $\partial_{q_l} \cdot q_m$ 提供不同圈动量之间的关系，是完备约化系统所必需的。

\textbf{验证}：独立 loop-scalar-products 数目 $N_{\text{sp}} = L(L+1)/2 + LK$ 与需要闭合的 $z/ISP$ 坐标维度匹配。这里不把顶点能量符号（如独立 \texttt{ke[i]}）算入标量积空间。

\subsection{ISP 处理}

\textbf{定义}：给定拓扑的传播子 $\{\xi_e^2\}$，所有标量积 $\{q_l \cdot q_m,\, q_l \cdot k_j\}$ 中不能表示为 $\xi_e^2$ 线性组合的，称为 ISP（Irreducible Scalar Products）。

\textbf{函数族扩展}：积分家族增加 ISP 指标：
\begin{equation}
  J[\{a_v\}, \{\text{pack}_e\}, \{n_{\text{isp}_j}\}]
\end{equation}
ISP 本身只指不可约标量积的 numerator，指标定义零点固定为 $0$：$n_{\text{isp}_j}>0$ 时被积函数乘入该坐标的 $n_{\text{isp}_j}$ 次幂，$n_{\text{isp}_j}=0$ 时不含该 ISP。用户若显式给出 $n_{\text{isp}_j}<0$，package 保留并生成，但该对象物理上是同一坐标槽的额外 denominator，不再称为 ISP。Association 中即使另给 \texttt{zeroPoint}，初始化 metadata 仍规范为 $0$。

对只选择 numerator 区域的 family，target-to-seed 反推仍须容纳该 family 全部 IBP generator 的 ISP shift span。package 不会为此把下界从 $0$ 自动降到负值；若用户实际选择的 target range 使某些 template 的 seed 域为空，只能由用户明确扩大上界。这种上界扩张不改变 ISP 零点，也不把负幂重新引入当前 family。只生成 general seeds 的 family 不进入 target-to-seed 反推。

ISP 是二次齐次的 numerator 坐标，但其指标 $n_{\mathrm{isp}}$ 不因此接受额外的奇偶筛选。parity constraints 只能来自该 family 已推导的 line-power/endpoint-state 不变量；single-massive sunrise 的 odd/odd 条件中两个 ISP 槽的系数均为零。

Parity capability 逐 line 读取初始化后的函数系统：massive line 只接受已证明闭合的 h/H preset，massless line 接受 exponential quotient。后者由式~\eqref{masslessDerivativeCanonicalTech} 连同动量幂 shift 保持
\begin{equation}
 p_e=b_e+n_{e,1}+n_{e,2}\pmod 2.
\end{equation}
因此 pure-massless family 可以使用用户显式给出的 root constraints，并沿用同一 sector contact transport；package 不自动猜 remainder。独立的 seed-level \texttt{massless\_sunrise\_bundle\_guard} benchmark 使用三条 routing
\begin{equation}
 \xi_1=l_3,\qquad \xi_2=k_{321},\qquad
 \xi_3=l_3-k_{321}-k_L,
\end{equation}
实际 general-template shifts 分别保持 $p_1,p_2,p_3$。该 seed-level benchmark 选择 odd/odd 子空间 $p_1=1$、$p_2+p_3=1$；两个 ISP 槽在这两行中的权重均为零。它不是另一个长期 sunrise 全流程 example，也不进入撒点或 reduction。未知或 custom function system 没有该证明时，显式 parity 请求仍 fail closed。

\textbf{用户输入}：
\begin{verbatim}
loopMomenta = {l3, k321};
loopExternalMomenta = {wdnmd};
independentExternalMomenta = {kE1};
vertexEnergies = <|1 -> ke[1], 2 -> Sqrt[sp[wdnmd,wdnmd]],
  3 -> Sqrt[sp[kE1,kE1]]|>;

ispData = {
  {rhoA, sp[l3, wdnmd + l3], {0, 1}},
  {rhoB, sp[k321, wdnmd], {0, 1}}
};
\end{verbatim}

\texttt{sp[p,r]} 是用户口 scalar-product convention，并具有 \texttt{Orderless} 属性。所有含圈 line momentum 与 ISP 参数必须由 \texttt{loopMomenta} 和 \texttt{loopExternalMomenta} 线性张成；实际无圈 line/phase 模长必须由 \texttt{independentExternalMomenta} 完备声明。017 对 loop Gram 使用短名 \texttt{ssij}，对独立无圈模长使用 \texttt{sE1,sE2,...}。类别总数超过 9 时按总数位宽补零。与两类动量声明都无关的相位标量仍可写成 \texttt{ke[i]}。

下文统一把 \texttt{loopExternalMomenta} 的按序元素简称为 $k_{L1},k_{L2},\ldots$（\texttt{kL}），把 \texttt{independentExternalMomenta} 的按序元素简称为 $k_{E1},k_{E2},\ldots$（\texttt{kE}）。简称只编码列表角色与位置，不约束用户符号名。两套编号独立：\texttt{kL} 进入内部 \texttt{qk/kk} 与公开 \texttt{ssij}；\texttt{kE} 的平方模长进入独立的 \texttt{externalLegSquaredCoordinate[e]} 与公开 \texttt{sEe}，不进入 loop scalar-product 空间。

\textbf{完备性验证}：\texttt{verifyISP[topology, ispData]} 检查：
\begin{enumerate}
  \item 所有标量积 $\{q_l \cdot q_m,\, q_l \cdot k_j\}$ 均可表示为 $\{\xi_e^2\}$ 和 $\{\text{isp}_j\}$ 的线性组合
  \item ISP 之间线性无关；当前主线中 ISP 表达式可为 \texttt{sp[p,r]} 或其线性组合坐标，不要求直接是某个内部编号变量
  \item \texttt{zExprs} 与 ISP 坐标总数等于独立 loop-scalar-products 数量，即 $\# z_e+\#\text{ISP}=N_{\text{sp}}$
  \item line momentum 与 \texttt{sp[p,r]} 参数必须是声明动量基的线性组合；非线性输入会触发 \texttt{nonLinearLineMomenta} 或 \texttt{nonLinearScalarProductArguments}
  \item 数量闭合后必须能实际反解出 \texttt{repSP2Z}；重复或退化传播子动量会触发 \texttt{scalarProductCoordinateSolveFailed}
  \item 数值规则按工作流分类。不再对外部变量求导的 numeric linear 工作流可为 \texttt{ssij/sEe}、自定义坐标或顶点能量给定值；但 Kira 结果若要进入 \texttt{DSDE}，全部 derivative variables、其内部平方原子以及任何含它们的规则右端都必须保持符号。此时只允许固定 \texttt{dim/epsilon}、\texttt{nu}、zero-point 和 normalization 等非微分系数参数。
\end{enumerate}

这一步验证用户初始化给出的 \texttt{z/ISP} 坐标系是否闭合。程序不把 dS 图默认理解为 overcomplete propagator family，也不自动从传播子中选择独立子集；若输入中存在重复/退化传播子、ISP 不足或特殊数值外动量导致不可反解，应由用户修正 family 输入。
验证通过后，才进入 IBP 生成步骤。

用户确认的离散图对称性通过 \texttt{symmetryRules} 输入，并由 \texttt{symmetry}/\texttt{repSymmetry0} 在 EOM 后的 seed canonicalization 中消费。用户必须先为每个等价类选定唯一 canonical representative，并只提供从非代表到代表的单向规则；不能同时提供 $A\to B$ 与 $B\to A$。package 当前不替用户给未定向等价关系排序，也不使用 \texttt{ReplaceRepeated}。未来可另加低优先级 helper，以积分复杂度和稳定字典序作为缺省排序键自动定向，但它不属于当前 general seeds/operators 门禁。对 single-massive sunrise，两个顶点的外腿能量都取 \texttt{kE}，圈外动量模长为 \texttt{kL}；顶点交换反转每条 full line 的两个 endpoint slots。两条 massless 平行线交换还必须同步交换 line 2/3 的 pack 和 ISP 1/2 的指数。为使后一步是严格置换，ISP 可选为
\begin{equation}
 \rho_2=k_{321}\!\cdot l_3,\qquad
 \rho_3=(l_3-k_{321}-k_L)\!\cdot l_3,
\end{equation}
其中顶点交换、line 2/3 交换及 ISP 1/2 指数交换必须由同一个用户对称性模块作用于统一三参数 \texttt{J[aList,linePacks,ispList]}，不能另建 sunrise 专用积分 Head。该 example 的 \texttt{DSSeeds} 按正式接口构造全部 contact-reachable sector 的 general 模板，连续指标始终保持符号；它同时读取初始化生成的 $\{\texttt{ss11},\texttt{kE}\}$ general 参数微分算符，其中
\begin{equation}
 \frac{\partial}{\partial\texttt{ss11}}
 =2\,\texttt{ss11}\frac{\partial}{\partial\,sp[k_L,k_L]},
\end{equation}
而共同参数 \texttt{kE} 的算符同时作用于两个顶点能量方向。该 example 和对应独立手推不调用 \texttt{DSMetaSeedRange}、\texttt{DSGenerateIBP}、\texttt{DSLinear}、Kira、DE 或 scaling。
它们在 $k_{321}\leftrightarrow l_3-k_{321}-k_L$ 下互换。只交换 line pack 而保留旧 ISP 指数不是该图的对称性，不得写入 canonical rule。

\subsection{Seed shift metadata 与目标包络反推}

017 把 ``general seed 中有哪些指标移位'' 与 ``用户最终要生成多大的整数域'' 分成两个步骤。首先调用
\begin{verbatim}
seedData = DSSeeds[context];
allSeeds = DSAllSeeds[seedData];
seedGroups = DSSeedGroups[seedData];
groupInfo = DSSeedGroupMetadata[seedData];
shiftMeta = DSMetaSeedRange[seedGroups,{a1,a2,b1,b2,...}];
\end{verbatim}

\texttt{DSSeeds} 完整遍历物理独立的离散基底：massive 端点为 $n_i=0,1$ 四态，masslessFull 为 $n_2=0$ 的 $00/10$ representatives；随后先执行 EOM、symmetry 与 canonical，连续指标保持 general。它同时返回：
\begin{itemize}
  \item \texttt{allSeeds}：保持稳定 ordinal 的一维密封模板列表；
  \item \texttt{seedGroups}：依次按 $\{\texttt{sectorKey},\texttt{ibpClass},\texttt{generator}\}$ 分组，每个二级列表对应一个时间 IBP 顶点或一个复合圈动量 IBP 算符，并包含该算符的全部离散态模板；
  \item \texttt{seedGroupMetadata}：与上述分组同序的 sector、IBP 类型、生成元、模板 ordinal 与计数。
\end{itemize}
默认情况下 \texttt{DSSeeds} 用推荐的 \texttt{seedGroups} 初始化一次 shift metadata。用户可随时再次调用 \texttt{DSMetaSeedRange} 覆盖此前状态：若输入是 flat 列表 $\{s_1,s_2,\ldots\}$，则整体统计为一组；若输入含二级或更深列表，则只以最外层元素分组，每个元素内部无论嵌套多深都先完全 \texttt{Flatten} 后独立统计。因此用户既可逐 IBP 算符合并全部 $n_i$ 状态，也可把每条离散模板分别包装成顶层组。调用时程序用中英双语报告该规则、实际组数和来源。

shift 提取不读取 \texttt{J[aList,linePacks,ispList]} 的槽位长度。对每个积分先把 \texttt{J} 的全部参数统一 \texttt{Flatten} 成指标数据，再用 \texttt{Variables} 得到实际变量集合。用户声明的指标列表只作为完整性审计：声明遗漏的变量记入 \texttt{missingIndices} 并自动补入，声明中不存在于 seed 的变量记入 \texttt{extraIndices} 且不参与撒点；两者均给 warning，但不终止 metadata 初始化。每个输入组、每个实际指标 $x$ 分别保存所有仿射出现
\begin{equation}
 x+\delta,\qquad \delta\in\mathbb Z,
\end{equation}
的 shift 集、\texttt{shiftBounds}$=\{\min\delta,\max\delta\}$ 以及相对目标包络的 offset。含多个指标的 contact 组合（例如 $a_1+a_2+c$）不用于单轴 shift 统计，而由源 sector 方程自然带入 child sector。

随后用户给最终关系允许出现的目标指标包络。统一形式
\begin{verbatim}
DSGenerateIBP[allSeeds,{L,U}]
\end{verbatim}
把同一整数区间用于全部 root 指标；逐指标形式
\begin{verbatim}
DSGenerateIBP[allSeeds,{a1,La1,Ua1},{b1,Lb1,Ub1},...]
\end{verbatim}
要求 root 指标 exact cover。shrink 指标 \texttt{bS[e]} 继承同一 root line 的 \texttt{b[e]} 包络；若两者都显式出现且区间相同则合并，若区间冲突则拒绝该范围输入。

对某组、某指标的全部 shift $\Delta_x$，seed 值必须同时满足
\begin{equation}
 L_x\leq x+\delta\leq U_x\qquad (\forall\delta\in\Delta_x).
\end{equation}
因此正确的反推区间是逆像的交集
\begin{equation}
 x\in\bigl[L_x-\min\Delta_x,\;U_x-\max\Delta_x\bigr].
 \label{seedEnvelopeInverse017}
\end{equation}
对普通指标使用式~\eqref{seedEnvelopeInverse017}。对 ISP 指标还要保持“package 不自行降低用户下界”的边界：
\begin{equation}
 x_{\rm ISP}\in\bigl[\max\{L_x,L_x-\min\Delta_x\},\;U_x-\max\Delta_x\bigr].
 \label{ispSeedEnvelope018}
\end{equation}
用户显式给出的 $L_x<0$ 仍完整有效；式~\eqref{ispSeedEnvelope018} 只禁止 planner 为了覆盖升幂项把下界进一步推向更负的整数。对所选 seed 作用全部 shift 后，所得关系仍落在用户 target envelope 内，但不再承诺为了填满下边界而自动生成更低 ISP seed。不得误用并集外包络 $[L_x-\max\Delta_x,U_x-\min\Delta_x]$，后者会无故扩大撒点域。若修正后的下界大于上界，该组为空；程序 warning 并保留结构化 \texttt{emptySeedRangeGroups}，但该结果的 complete IBP capability 为 false，不能作为完整后端约化系统。

实际生成顺序固定为：按式~\eqref{seedEnvelopeInverse017} 构造候选整数点，\emph{在生成元作用前}按当前 sector 的 parity 筛点，再把连续指标值代入 general seed，应用用户已经单向定向的 sector symmetry，随后重复代入同一组精确 numeric rules 并对每个线性项执行 \texttt{Together/Cancel}。\texttt{DSSeeds[...,ApplyNumericRules->True]} 在 general template 密封前已经做第一次物理规则代入；此时 \texttt{a/b/bS/ispN} 仍是合法 continuous indices，必须与其它 \texttt{seedResidualCoefficientVariables} 分开。撒点后 \texttt{sampledCoefficientVariables} 才表示仍未数值化的真实 coefficient residual。纯数值工作流必须先按符号 topology 构造参量微分算符，再把同一个固定非奇异精确有理点写入 seed 的 \texttt{numericRules}；验收要求 \texttt{seedResidualCoefficientVariables} 与 \texttt{sampledCoefficientVariables} 都为空。逐关系的直接复核应收集全部不同的 \texttt{J}，提取各 \texttt{Coefficient} 和常数项，先精确重构原关系，再确认这些系数的 \texttt{Variables} 为空；不能把多积分关系除以任意一个 \texttt{J}。生成后的 parity certificate 只检查关系，不删除关系中的积分，也不把不满足 parity 的积分替换为零。\texttt{candidatePointCount}、\texttt{parityAcceptedPointCount}、逐组 shift/range 和空域状态均写入返回数据。用户包络唯一决定规模；package 不设置 seed 数、离散态数或方程数上限。

018 不改变上述代数顺序，但把连续撒点热路径实现为两层复用。一次 \texttt{DSGenerateIBP} 先由初始化的完整 sector metadata 构造
\begin{equation}
 \texttt{sectorKey}\longmapsto\{\texttt{sectorMetadata},\texttt{sectorTopology}\}
\end{equation}
缓存；每个数值积分仍由 root-ordered line-pack pattern 导出 sector key，并再次通过该 metadata 的 shape 匹配，故缓存不能绕过 sector gate。每条 general template 只在枚举前投影一次当前公开坐标；数值代入和 sector canonical 后，只有结果确实重新含有 \texttt{qq/qk/kk/z/xi/rho} 或无圈模长内部原子时才再次投影。该优化不缓存 canonical 结果，也不改变 EOM、symmetry、parity 或系数。

密封的 \texttt{DSSeeds} Association 模板保存 source integral、sector、generator、逐模板 \texttt{templateHash} 和全集 hash，因而继续在代入 general seed 前执行 parity 筛点。完整列表与合法子集分别标记为 \texttt{sealedComplete} 和 \texttt{sealedSubset}；子集不能继承全集的 complete 能力。缺省 \texttt{standard} 路线只核对 producer 写入的 set hash、计数和 ordinal metadata，不逐条重算模板 hash；显式 \texttt{full} audit 才重新计算逐模板与完整集合内容 hash。公开接口也允许直接给出只含 \texttt{J} 的 raw expression 列表；这类输入没有足够 provenance 来判断 source-seed parity，程序不得猜测，故保留用户点域并执行完整 consumer 审计。三条路线都可生成 backend-neutral partial linear data，但只有完整密封路线具有 \texttt{completeSystemQ=True}。

\texttt{DSGenerateIBP} 把 \texttt{sealedQ/sourceDigest/completeSystemQ/subsetQ/auditLevel/capabilities} 写入 \texttt{artifactContract}。缺省 \texttt{AuditLevel->"standard"} 依赖同一 producer 已完成的 parity 预筛与 sector canonical；\texttt{"full"} 用于 release/audit，重跑关系级 parity certificate。\texttt{DSLinear} 的 \texttt{standard} 路线只核对 sealed 状态、方程计数和 producer 的 \texttt{sourceDigest} 字段，不重新 hash 全部关系，不再次对全部关系执行 symmetry 或 coverage classifier；\texttt{full} 才重算 source digest 与 coverage，unsealed/raw 输入仍走完整 consumer 扫描。linearization 后按实际数学方程 \{coefficient rules, constant term, nonlinear terms\} 删除重复项，不把 generator/sector 等来源 metadata 放进相等键；返回 \texttt{sourceEquationCount}、\texttt{equationCount} 和 \texttt{duplicateEquationCount}。linearization 本身产生的方程、映射和 targets 在 Kira export 边界由 artifact content digest 覆盖，不另加重复的内部 \texttt{linearDigest}。partial 结果显式保持 \texttt{completeSystemQ=False}，不能通过 formal reduction 门禁。

在枚举前，\texttt{DSGenerateIBP} 逐输入组用中英双语打印反推范围，显示为
\begin{verbatim}
编号 1（来源）：{a1,min,max}, {b1,min,max}, ...
Group 1 (source): {a1,min,max}, {b1,min,max}, ...
\end{verbatim}
编号只表示用户输入的 seed 分组；即使缺省分组通常对应一个物理 IBP 生成元，也不把用户自定义分组称为 IBP 算符。来源取该组的 sector、IBP 类型与生成元；flat 总组则列出其中全部来源。若任一规则出现 $\min>\max$，程序另行 warning 并列出相同编号、来源和空区间。warning 明确解释：目标包络窄于该组 shift 跨度，故不存在能使全部移位后积分仍落在目标包络内的 seed 点；该编号的 IBP 撒点为空，最终方程集合缺少对应关系，不能作为完整约化系统。该组同时写入 \texttt{emptySeedRangeGroups}。全局消息开关可关闭显示，但不改变 metadata、方程集合或 capability 状态。

\subsection{分 Sector IBP 生成流程}

\begin{enumerate}
  \item 读取拓扑 + ISP 配置
  \item 验证 ISP 完备性
  \item 构造 IBP 生成元集合 $\{\mathcal{O}_{l,v}\}$（$L(L+K)$ 个，$K=\#\texttt{loopExternalMomenta}$）
  \item 对每个 sector（由缩并线集合标记）：
  \begin{enumerate}
    \item 构造该 sector 的连续指标盒子 $\{a_v, b_e, n_{\text{isp}}\}$ 与离散 $n$ 状态集合
    \item 按式~\eqref{seedEnvelopeInverse017} 从最终关系包络逐输入组反推连续 seed，并在生成元作用前筛选 parity
    \item 对每个输入组使用 \texttt{DSSeeds} 已完整枚举的独立离散模板；massive 为四态，masslessFull 为 $00/10$ quotient representatives；生产路径不允许离散抽样，也不允许残留符号 $n$ 或 massless source $n_2=1$
    \item 对每个离散态，作用该组对应的 time-IBP 或 momentum-IBP 生成元：
      \begin{itemize}
        \item 应用链式法则 $\to$ $z/\xi$-导数 + ISP-导数 + 相位项
        \item 转化为指标移位算符，包含传播子幂次项与 building-block 导数项
        \item 立即应用 EOM 递推，递归消去所有 $n\geq 2$
        \item 应用 massless 双 theta quotient relations，保持 $\{b_e,n_{e,1},n_{e,2}\}$ 指标包
        \item 扫描输出，若仍含 $n=2$ 或未知 pack，则该 seed 失败，不进入后端导出
      \end{itemize}
    \item 保存 EOM-canonical IBP 方程，命名区分 sector 与生成元类型
  \end{enumerate}
  \item 输出：按 sector 分组的 canonical IBP seed 文件。linear/Kira 阶段只能读取这些文件。
\end{enumerate}

\textbf{命名规则}：\texttt{IBP\_sector\_<shrunkLines>\_seed\_<seedIndex>.dat}

例如：\texttt{IBP\_sector\_none\_seed\_001.dat}（top sector），\texttt{IBP\_sector\_e3\_seed\_012.dat}（线 3 缩并）。

\subsection{后端中立的 linear-system 与 serializer}

canonical seed 先由 \texttt{makeLinearSystemData} 转成后端中立的 \texttt{linearData}。这一层不属于 Kira，也不包含任何后端可执行文件路径。主要数据为：
\begin{itemize}
  \item \texttt{linearEquations}：按全局积分编号表示的线性方程；
  \item \texttt{integralList} 与 \texttt{integralRules}：全 sector 统一排序的积分表及映射；
  \item \texttt{sectorMetadataList}：每个 sector 的 compact $a$ 槽、line pack、缩并和原拓扑对应；
  \item coefficient、ordering、coverage 与 readiness reports。
  \item \texttt{artifactContract}：producer source digest、linear digest、complete/subset 状态和 serializer/formal-reduction capability。
\end{itemize}

serializer 只把 \texttt{linearData} 转换为目标后端要求的语法、编号和文件布局，并执行后端特有的输入检查；它不重新生成 IBP seed，不重新应用 EOM，也不改变 sector convention。当前 Kira serializer 生成 \texttt{userSystem/ibp.kira}、\texttt{list}、\texttt{jobs.yaml}、积分映射和 metadata。未来若接 Rational Tracer，应新增独立 serializer，但 canonical seed 和 linear-system 主线无需重写。

\texttt{DSKiraPlan} 的 \texttt{preReduction} 阶段允许 partial system，并把它明确保留为不完整预分析；\texttt{formal} 阶段要求 \texttt{completeSystemQ=True}。formal plan 只构造一次 active-basis 一阶导数 closure，把完整结果保存为 \texttt{preparedLinearData}，并用 \texttt{analyticDerivativeCertificate} 固定其 hash；plan 驱动 export 必须校验证书并直接消费这份 payload，不得从原始 \texttt{linearData} 重建第二次。

初次 pre-reduction 探测不得把完整积分表设为 targets。候选数量按预估 master 规模设置；没有更具体 family 依据时取约 1000 的保守上限。formal 阶段只选择实际完整 active-basis token 与其解析一阶导数闭包。任意小批候选及其一阶导数只能证明局部覆盖，不能确定 master 数；必须先读取 pre-reduction 的实际 master 列表，再构造 formal targets。长期 single-massive sunrise example 在 general seeds/operators 阶段结束，不进入本段流程。

export manifest 的 \texttt{artifactIdentity} 同时固定 linear source/linear digest、实际 linear equations、积分映射、targets、coefficient rules、active-basis payload，以及写出的 \texttt{ibp.kira/list/jobs.yaml/maps/metadata} 文件 SHA-256。importer 逐项复算表达式与磁盘文件 digest；任一 map、target、rule、active payload 或导出文件变化均拒绝。package patch 版本字符串只作诊断，不再代替内容身份；同源性仍要求 context input hash、zero point、symmetry、tadpole 与 loop--tree convention 相同。

无论系数是否仍含实符号参数，Kira 输入都必须先实数化。所有 massless propagator 动量原子由 topology/line 的 phase dependency metadata 结构性识别，并只在 backend 执行 $k\mapsto-i\,ik$；不得按符号名称猜测。massive bubble 的既有实现作为参考路线。设 $x$ 表示全部实 backend 变量，则每个非零系数写成
\begin{equation}
  c_{ri}=q_{ri} I^{\sigma_{ri}},\qquad
  q_{ri}\in\mathbb{Q}(x)\setminus\{0\},\quad \sigma_{ri}\in\{0,1\}.
\end{equation}
维护侧最小检查必须同时覆盖符号实有理函数正例、全数值正例和不可分离实虚部负例，并直接审计生成的 \texttt{ibp.kira} 文本不含 \texttt{I}、\texttt{Complex} 或 \texttt{dsii}。该检查只验证 serializer convention；具体 family 仍须用其 formal targets 重新导出并在外部 Kira 中验收。

exporter 在完整积分 ID 图上求二元列相位 $p_i$，使同一方程内
\begin{equation}
  \sigma_{ri}+p_i=\rho_r\pmod 2.
\end{equation}
随后定义物理积分与 backend 积分的关系
\begin{equation}
  J_i=I^{p_i}K_i
\end{equation}
并把第 $r$ 行整体乘以 $I^{-\rho_r}$。所得 Kira 系数严格属于实有理函数域 $\mathbb{Q}(x)$；全数值路线进一步退化为 $\mathbb{Q}$。相位约束冲突、系数同时含非零实部和虚部、machine real 或其它非精确数值都必须在写文件前拒绝。把 $I$ 改名为 \texttt{dsii} 不构成实数化；磁盘上的 \texttt{ibp.kira} 出现 \texttt{I}、\texttt{Complex}、\texttt{dsii} 或其它虚数替代 token 时，export 必须失败。

\texttt{gaussianPhaseGauge} manifest 保存全部 ID 的 $p_i$、逐行 $\rho_r$、约束数、连通分量、冲突数和 convention。若 Kira 返回
\begin{equation}
  K_a=\sum_b R_{ab}K_b,
\end{equation}
importer 必须先恢复
\begin{equation}
  J_a=\sum_b I^{p_a-p_b}R_{ab}J_b,
\end{equation}
再应用 ID 到 \texttt{J} 的映射。manifest 缺项、相位不是 $0/1$ 或 ID 未完整覆盖时拒绝导入。

Kira~2.3 的 FireFly 路线要求至少一个 coefficient symbol；零变量系统会报 \texttt{No variables?}。因此纯有理 user-system 不追加 \texttt{ccc} dummy，也不写 \texttt{dsii}，而使用 Kira 官方的 \texttt{run\_triangular:true} 与 \texttt{run\_back\_substitution:true} 常系数路线，并关闭 \texttt{run\_firefly}。prepare 必须同时检查前端 coefficient variables、backend variables、dummy 状态和磁盘上的实际 \texttt{ibp.kira} 文本。

积分 ID 的唯一顺序来源是 \texttt{linearData["integralList"]}。\texttt{DSKiraPlan}、active-basis preparation、公开 exporter 和底层 serializer 都只能读取该顺序，不得按 preferred master、complexity 或 sector 再次编号。用户确需改变顺序时，必须在 backend 边界前显式调用 \texttt{DSReorderIntegrals[linearData,order]}；该函数一次性同步更新 \texttt{integralList}、\texttt{integralRules}、线性方程 ID 与顺序 digest。\texttt{preferredIntegrals} 只表示候选偏好，不改变全局 ID；旧 \texttt{KiraIntegralOrder} 不再是 export option。这样 manifest、双向映射、reduction 与 DE 使用同一份冻结编号，避免 plan 与 serializer 分别重排造成同一个 ID 指向不同 \texttt{J}。

本 package 默认不安装、配置或运行 Kira/Rational Tracer，也不保存本机 Kira/Fermat 路径。018 的 importer 可读取并验证用户在 package 外生成的完整 reduction 与 master-integral 结果。实际方程/文件内容 identity、双向积分映射、targets、master 顺序、Gaussian phase/energy convention、系数域及 reduction closure 是硬边界；\texttt{kira.log} 的完成标记只作为诊断 warning，因为日志措辞或缺失不能推翻已经通过的结构化闭合证据。

% ============================================================
\section{标量积变量与 $z=\xi^2$ 线性变换}
% ============================================================

本节引入 $z_e = \xi_e^2$ 变量，建立标量积到 $z$ 变量的线性变换体系。这一框架是动量 IBP 链式法则中指标移位操作的基础。

\subsection{$z$ 变量定义}

对每条内线 $e$，定义
\begin{equation}
  z_e \equiv \xi_e^2 = Q_e \cdot Q_e,
  \label{zDef}
\end{equation}
其中 $Q_e^\mu = \sum_l c_{e,l}\, q_l^\mu + P_e^\mu$ 为线 $e$ 的总动量。$z_e$ 与 $\xi_e$ 的关系为
\begin{equation}
  \xi_e = z_e^{1/2}.
  \label{xiFromZ}
\end{equation}

传播子的物理幂次用 $z$ 变量表示为
\begin{equation}
  \xi_e^{-(b_e + b0_e)} = z_e^{-(b_e + b0_e)/2}.
  \label{zPower}
\end{equation}
因此 $z_e^n$ 作用在 $z_e^{-(b_e+b0_e)/2}$ 上等价于指标移位 $b_e \to b_e - 2n$。

\subsection{标量积变量约定}

独立标量积分为三类，其中 $K=\#\texttt{loopExternalMomenta}$ 只统计进入 loop scalar-product 空间的外部三动量向量：

\begin{itemize}
  \item \textbf{圈--圈}：$q_l \cdot q_m$（$l \leq m$），共 $L(L+1)/2$ 个。
  \item \textbf{圈--外}：$q_l \cdot k_j$，共 $LK$ 个。
  \item \textbf{外--外}：$k_i \cdot k_j$ 定义为独立符号，不保持矢量点积结构；017 未指定 \texttt{KinematicRules} 时，按 \texttt{loopExternalMomenta} 的位置生成 $\operatorname{sp}(k_i,k_j)\to ss_{ij}^{,2}$。在 IBP 生成中这些量作为常数参数处理。
\end{itemize}

总独立标量积数目为
\begin{equation}
  N_{\text{sp}} = \frac{L(L+1)}{2} + LK.
  \label{Nsp}
\end{equation}

将所有独立标量积排成向量 $\vec{s}$，其分量依次为 $\{q_1^2,\, q_1 \cdot q_2,\, \ldots,\, q_L^2,\, q_1 \cdot k_1,\, \ldots\}$。顶点相位中的 $|k|$ 或 $|k_a|+|k_b|$ 是能量参数，不是 $\vec{s}$ 的分量。

\subsection{线性变换的推导}

线 $e$ 的总动量平方为
\begin{equation}
  z_e = Q_e^2 = \left(\sum_l c_{e,l}\, q_l + P_e\right)^2.
  \label{zeExpand}
\end{equation}

展开后：
\begin{equation}
  z_e = \sum_{l,m} c_{e,l}\, c_{e,m}\, (q_l \cdot q_m) + 2 \sum_l c_{e,l}\, (q_l \cdot P_e) + P_e^2.
  \label{zeExpandFull}
\end{equation}

其中 $P_e$ 为外动量组合，$P_e^2$ 和 $q_l \cdot P_e$ 中的外动量点积均以当前动力学坐标替代；018 缺省写成 $ss_{ij}^{,2}$，也可由用户重定义。因此 $z_e$ 是标量积变量 $\vec{s}$ 的线性函数加上常数项。

对所有线 $e$ 写为矩阵形式：
\begin{equation}
  \vec{z} = M \vec{s} + \vec{c},
  \label{zLinear}
\end{equation}
其中 $M$ 是 $N_z \times N_{\text{sp}}$ 矩阵（$N_z$ 为用户输入的传播子平方方程数），$\vec{c}$ 的分量为各 $P_e^2$（用缺省 $ss_{ij}^{,2}$ 或用户自定义坐标表达式表示）。

当 \texttt{zExprs} 数等于独立标量积数（$N_z = N_{\text{sp}}$，即无 ISP）时，$M$ 为方阵且可逆，逆变换为
\begin{equation}
  \vec{s} = M^{-1}(\vec{z} - \vec{c}).
  \label{zInverse}
\end{equation}

当存在 ISP 时，逆变换需补充 ISP 的定义关系（ISP 直接用 $\vec{s}$ 中的分量表示）。当前实现要求补充后得到方阵闭合系统；若用户给出的传播子平方方程相对 ISP 过多或不足，validation 报告会要求修正输入，而不是自动挑选子集。

\subsection{Bubble 拓扑的显式计算}

以 bubble 拓扑（单圈，$L=1$，$E=2$ 条内线，1 个外动量 $k$）为例。两条内线的动量为
\begin{align}
  Q_1 &= q_1, \\
  Q_2 &= q_1 - k.
\end{align}

\textbf{正变换} $\vec{z} = M \vec{s} + \vec{c}$：
\begin{align}
  z_1 &= Q_1^2 = q_1^2, \label{z1bubble} \\
  z_2 &= Q_2^2 = q_1^2 - 2\, q_1 \cdot k + s_{11}, \label{z2bubble}
\end{align}
其中 $s_{11}$ 是默认外部不变量名，对应 $k \cdot k$；若用户自定义则替换为对应变量名。矩阵形式为
\begin{equation}
  \begin{pmatrix} z_1 \\ z_2 \end{pmatrix}
  = \begin{pmatrix} 1 & 0 \\ 1 & -2 \end{pmatrix}
  \begin{pmatrix} q_1^2 \\ q_1 \cdot k \end{pmatrix}
  + \begin{pmatrix} 0 \\ s_{11} \end{pmatrix}.
  \label{bubbleMatrix}
\end{equation}

\textbf{逆变换} $\vec{s} = M^{-1}(\vec{z} - \vec{c})$：
\begin{align}
  q_1^2 &= z_1, \label{q1sqInv} \\
  q_1 \cdot k &= \frac{1}{2}(z_1 + s_{11} - z_2). \label{q1kInv}
\end{align}

\textbf{标量积 $q_1 \cdot Q_e$}（在 IBP 链式法则中直接使用）：
\begin{align}
  q_1 \cdot Q_1 &= q_1^2 = z_1, \label{q1Q1} \\
  q_1 \cdot Q_2 &= q_1 \cdot (q_1 - k) = q_1^2 - q_1 \cdot k = z_1 - \frac{1}{2}(z_1 + s_{11} - z_2) = \frac{1}{2}(z_1 + z_2 - s_{11}). \label{q1Q2}
\end{align}

这些关系表明 $q_l \cdot Q_e$ 可以表示为 $z$ 变量和外动量不变量（缺省 $ss_{ij}^{,2}$ 或用户自定义坐标表达式）的线性组合。

\subsection{在动量 IBP 链式法则中的应用}

动量 IBP 的核心算符为 $q_l^\mu\, \partial / \partial q_l^\mu$（对角生成元）或 $q_m^\mu\, \partial / \partial q_l^\mu$（交叉生成元）。通过链式法则将 $\partial / \partial q_l^\mu$ 分解为对 $\xi_e$ 的导数：
\begin{equation}
  q_l \cdot \frac{\partial}{\partial q_l} = \sum_e c_{e,l}\, \frac{q_l \cdot Q_e}{\xi_e}\, \frac{\partial}{\partial \xi_e}.
  \label{chainRuleXi}
\end{equation}

利用 $\xi_e=z_e^{1/2}$，有
\begin{equation}
 \frac{\partial}{\partial\xi_e}
 =2\xi_e\frac{\partial}{\partial z_e},
 \qquad
 z_e\frac{\partial}{\partial z_e}
 =\frac{\partial}{\partial\ln z_e}.
\end{equation}
因此
\begin{equation}
  q_l \cdot \frac{\partial}{\partial q_l}
  =2\sum_e c_{e,l}(q_l\cdot Q_e)\frac{\partial}{\partial z_e}
  =2\sum_e c_{e,l}\frac{q_l\cdot Q_e}{z_e}
     \frac{\partial}{\partial\ln z_e}.
  \label{chainRuleZ}
\end{equation}

\textbf{指标移位规则}：$q_l \cdot Q_e$ 用式~\eqref{zInverse} 或显式计算表示为 $z$ 变量的线性组合后，每一项形如 $\alpha\, z_e$（$\alpha$ 为常数或外部不变量名的函数）。$z_e$ 作用在传播子 $z_e^{-(b_e+b0_e)/2}$ 上的效果为
\begin{equation}
  z_e \cdot z_e^{-(b_e+b0_e)/2} = z_e^{-(b_e - 2 + b0_e)/2},
  \label{zShift}
\end{equation}
即指标移位 $b_e \to b_e - 2$。

类似地，$1/z_e$ 项产生 $b_e \to b_e + 2$ 的移位。

\textbf{实现步骤}：
\begin{enumerate}
  \item 用替换规则将每个 $q_l \cdot Q_e$ 表示为 $\{z_e\}$ 和外部不变量（缺省 $ss_{ij}^{,2}$ 或用户自定义坐标表达式）的线性组合。
  \item $2\,\partial/\partial z_e$ 作用在 $z_e^{-(b_e+b0_e)/2}$ 上，产生系数 $-(b_e+b0_e)$ 和基础移位 $b_e\to b_e+2$。
  \item 再把 $q_l\cdot Q_e$ 中的每个 $z_e^n$ 或 ISP 单项式吸收到指标；$z_e^n$ 对应进一步移位 $b_e\to b_e-2n$。
\end{enumerate}

这一框架使得动量 IBP 的链式法则完全转化为指标移位操作，与式~\eqref{zDef}--\eqref{zShift} 定义的指标体系自洽。

% ============================================================
\section{独立变量微分方程求导}
% ============================================================

微分方程 seed 需要把 $\partial_x J$ 写回同一积分族。017 将独立变量分为三类：由 \texttt{loopExternalMomenta} 完整 Gram 坐标得到的模长、\texttt{independentExternalMomenta} 声明且在 topology 中实际出现的无圈动量模长，以及与任何动量向量都无关的独立标量能量。第一类沿用标准外动量矢量导数；第二类对绑定的 line momentum 作标量径向导数并同时微分相位；第三类只按显式系数与顶点相位依赖求导。构造 $D_{ij}$ 时，$i,j$ 只遍历 \texttt{loopExternalMomenta}。

\subsection{顶点能量参数}

设顶点相位写成
\begin{equation}
  v=+:\quad \exp[-iE_v\tau_v],
  \qquad
  v=-:\quad \exp[+iE_v\tau_v],
\end{equation}
并记
\begin{equation}
  \eta_v=
  \begin{cases}
  -i,& v=+,\\
  +i,& v=-.
  \end{cases}
\end{equation}
若 $y$ 是独立顶点能量参数，例如 \verb|ke[i]|，则它不参与动量标量积空间，只给出
\begin{equation}
  {\partial\over\partial y}J
  =
  \sum_v
  \left[-\eta_v\,{\partial E_v\over\partial y}\right]
  J[a_v\to a_v+1].
  \label{vertexEnergyDerivative}
\end{equation}
式中负号来自时间幂 convention：$J$ 中保存的是 $(-\tau_v)^{a_v+a0_v}$，而相位求导产生的是 $\eta_v\tau_v=-\eta_v(-\tau_v)$。

\subsection{平方 Gram 原子与根号坐标的矢量导数分解}

先令内部平方 Gram 原子为
\begin{equation}
  x_a\in\{k_i\cdot k_j\}_{1\le i\le j\le K},
\end{equation}
并定义外动量矢量导数算符
\begin{equation}
  D_{ij}=k_i\cdot{\partial\over\partial k_j}.
\end{equation}
它在 Gram 坐标上的作用为
\begin{equation}
  D_{ij}(k_m\cdot k_n)
  =
  \delta_{jm}(k_i\cdot k_n)
  +
  \delta_{jn}(k_i\cdot k_m).
  \label{externalVectorGramAction}
\end{equation}

因此先在约束坐标 $\{x_b\}$ 上解平方原子导数
\begin{equation}
  {\partial\over\partial x_a}
  =
  \sum_{ij}c^{(a)}_{ij}D_{ij},
  \qquad
  \sum_{ij}c^{(a)}_{ij}D_{ij}x_b=\delta_{ab}.
  \label{externalInvariantDecomposition}
\end{equation}
完整 $K^2$ 个 $D_{ij}$ 通常有零空间，所以 $c^{(a)}_{ij}$ 不唯一。原子实现选择上三角 external-vector basis 给出一组 canonical 解。017 不重写该分解；对直接模长
\begin{equation}
  r_a=\sqrt{x_a}=\sqrt{\operatorname{sp}(k_i,k_j)}
\end{equation}
只在外层应用
\begin{equation}
  {\partial\over\partial r_a}=2r_a{\partial\over\partial x_a}.
  \label{rootInvariantChainRule}
\end{equation}
显式旧平方坐标保持单位 Jacobian。实际返回接口是
\begin{center}
  \texttt{makeExternalInvariantDerivativeDecomposition},
\end{center}
它返回矩阵、所选系数、残差、\verb|nullity| 与 \verb|nonUniqueQ|。若某一物理问题要求不同切向代表元，应替换 operator basis，而不是把任意选择隐藏在求导核心中。

每个 $D_{ij}$ 作用到被积函数时，传播子、massive/massless building block 与 ISP 的求导都按同一链式法则处理。若某个顶点能量写成 $\sqrt{\operatorname{sp}(k_i,k_j)}$，则同一原子分解与式~\eqref{rootInvariantChainRule} 也作用到相位能量。若写成 $\sqrt{\operatorname{sp}(k_{E,i},k_{E,j})}$，它不产生 $D_{ij}$ 方向，只按显式标量依赖求导。

对一般非线性函数 $F(x_1,\ldots,x_N)$，程序必须显式使用
\begin{equation}
  D_{ij}F
  =\sum_a {\partial F\over\partial x_a}D_{ij}x_a.
  \label{externalVectorFunctionChainRule}
\end{equation}
因此 $F=\sqrt{\operatorname{sp}(k_1,k_1)}$ 给出普通链式因子。把每个坐标直接替换成 $D_{ij}x_a$ 只对线性 $F$ 成立，不能用于顶点能量或其它非线性动力学量系数；实现也不得用 \texttt{PowerExpand} 假定根号分支。

对 massive building block，loop-momentum IBP、time-IBP 与这里的 external-vector/外不变量导数共享同一个 line-local \texttt{compiledFunction\allowbreak System}：普通导数只读取最终 $A_T$ 编译的 \texttt{derivative\allowbreak Terms}。最终 $W_T=\det(T)W$ 及其 \texttt{shrink\allowbreak Terms} 只在 time-IBP 的 theta coincidence/Wronskian shrink 中使用；动力学量导数不微分 theta ordering，因此不调用 $W_T$。

\subsection{实际无圈动量线的径向导数}

设不含圈动量的 line momentum 为 $P_{E,e}$，并把其模长绑定为
\begin{equation}
  r_e=|P_{E,e}|=\sqrt{\operatorname{sp}(P_{E,e},P_{E,e})}.
\end{equation}
它仍是一条对 $\tau$ 积分的传播子，只是不产生新的 loop IBP generator。记该线物理分母幂为 $B_e=b_e+b0_e$。若端点 $s$ 的已编译基底满足
\begin{equation}
  {\partial f_{\alpha_s}(x)\over\partial x}
  =\sum_{\beta,p}c^{(e,s)}_{\alpha_s\beta p}\,x^p f_\beta(x),
  \qquad x=-r_e\tau_{v_s},
\end{equation}
则 $\partial_{r_e}$ 在指标表示中的 line-local 作用为
\begin{align}
 {\partial J\over\partial r_e}\bigg|_{e}
 ={}&-B_e\,J[b_e\to b_e+1]\nonumber\\
 &+\sum_{s=1}^{2}\sum_{\beta,p}
 c^{(e,s)}_{\alpha_s\beta p}\,
 J\!\left[
 a_{v_s}\to a_{v_s}+p+1,
 b_e\to b_e-p,
 n_{e,s}:\alpha_s\to\beta
 \right].
 \label{externalLegRadialDerivative}
\end{align}
这正是 \texttt{AT} 编译出的每一条 \texttt{derivativeTerms} 对 $a/b/n$ 的同时移位；其它端点态与其它线指标保持不变。对 h 基底的 $n=0$，$\partial_xh_0=h_1$，式~\eqref{externalLegRadialDerivative} 退化为两端各一个 $a_{v_s}\to a_{v_s}+1,n_{e,s}:0\to1$，再加分母项 $-B_eJ[b_e+1]$。

massless full/cross 线使用各自已固定的端点 phase 符号：full 线按实际求导端点翻转对应槽后立即应用 quotient canonical，cross 线保持固定的两个零槽；两者都加同一个分母幂导数。径向求导不微分 theta ordering，不读取 \texttt{WT/shrinkTerms}，所以不会把 time-IBP 的共同-theta contact 规则误用于微分方程。

若 $r_e$ 还进入顶点能量，则另加式~\eqref{vertexEnergyDerivative} 的相位项。对用户变量 $u$，实际组合为
\begin{equation}
 {\partial J\over\partial u}
 =\sum_{a\in\mathrm{loop\ Gram}}
 {\partial x_a\over\partial u}{\partial J\over\partial x_a}
 +\sum_{e\in\mathrm{no\!\mbox{-}\!loop\ magnitudes}}
 {\partial r_e\over\partial u}{\partial J\over\partial r_e}
 +{\partial J\over\partial u}\bigg|_{\mathrm{explicit\ scalar\ phase}}.
 \label{generalKinematicCoordinateChainRule}
\end{equation}
第一项的平方 Gram 原子导数已经包含由 $x_a$ 引起的顶点相位变化；最后一项只微分内部表达式中仍显式依赖 $u$ 的无圈模长或独立标量，不能重复加入 loop phase。

\subsection{变量重选、秩与零空间门禁}

把完整基础平方坐标按固定顺序写成
\begin{equation}
 \bm y=(\{x_{ij}\}_{i\le j},\{r_e^2\})^T,
\end{equation}
其中第二组是 topology 实际出现的无圈模长平方在第一组及此前已选无圈行之上的增量独立基；不把外腿向量点积作为公开坐标。每个未入基候选仍保存为 $\bm y$ 的线性组合，并在 line、phase 和显式系数中以该组合的平方根出现。用户规则统一写成
\begin{equation}
 A\bm y=\bm f(\bm u).
 \label{kinematicRuleMatrix}
\end{equation}
左端必须满足 $\operatorname{rank}A=N_y$。若不满秩，$\operatorname{NullSpace}(A)$ 给出未覆盖的基础动力学方向；$\operatorname{NullSpace}(A^T)$ 给出规则行依赖，其与 $\bm f$ 的收缩给出一致性残差。选取独立行解出 $\bm y(\bm u)$ 后，再构造
\begin{equation}
 B_{a\alpha}={\partial y_a(\bm u)\over\partial u_\alpha}.
\end{equation}
还必须有 $\operatorname{rank}B=N_y$。$\operatorname{NullSpace}(B^T)$ 按同一 \texttt{baseCoordinateOrder} 给出被用户参数化施加的局部约束方向；$\operatorname{NullSpace}(B)$ 给出冗余用户参数方向。

公开的两阶段调用为
\begin{center}
 \texttt{DSKinematics[input]},\qquad
 \texttt{DSKinematics[input,rules]},\\
 \texttt{DSInit[input,KinematicRules->rules]}.
\end{center}
第一条同时返回 topology/routing、显式双动量列表、缺省规则、从属模长 binding、可复制修改的 \texttt{selectionTemplate} 和上述秩/零空间；第二条只审计候选规则，第三条才初始化。该接口不会用 \texttt{Input[]} 阻塞 notebook、命令行或 batch。欠完备声明或规则是 error：返回 missing/null-space 证据并令初始化失败，后续 seed、求导、linearData 与 serializer 都由 capability gate 拒绝。过完备允许初始化并生成 symbolic IBP seed，但 $\texttt{derivativeUsableQ=False}$，故 \texttt{ds/DSDE} 与唯一 \texttt{rep2innerform} 均禁用。只有 exact 且满秩的坐标系统可按式~\eqref{generalKinematicCoordinateChainRule} 求导。

\subsection{带动力学量系数的表达式总导数}

微分方程基或约化结果一般是
\begin{equation}
  \mathcal I(s)=\sum_r c_r(s)J_r(s)+c_0(s),
\end{equation}
而不是单个指标积分。公开接口 \verb|ds[expr,var]| 使用 family 初始化时的外部变量名；018 缺省为根号坐标的 API 短名，并严格实现
\begin{equation}
  {\partial\mathcal I\over\partial s}
  =\sum_r c'_r(s)J_r(s)
  +\sum_r c_r(s){\partial J_r\over\partial s}
  +c'_0(s).
  \label{expressionTotalDerivative}
\end{equation}
实现时先把各个 $J_r$ 替换为惰性 token，让 Mathematica 的 \texttt{D} 只作用于显式系数；随后逐项调用单积分导数核心，合并后再统一执行 EOM、symmetry 与 canonical。接口拒绝内部 \texttt{kk[i,j]} 名称及 $J_iJ_j$ 等非线性输入，输出恢复到根号坐标或用户显式声明的兼容表示。

历史 012 阶段曾对十个 family 得到 $468/468$ 的 general-index 总导数比较，并对旧 reference bubble 得到 $80/80$ 的单积分与总导数比较；这些数字只作为旧版本证据保留，不代表当前 018 独立 benchmark 已重跑。当前任务书的 Phase~1 只在固定分支上手推并比较 general IBP seeds 与 general 参数微分算符，其中 single-massive sunrise 到此结束；Phase~2 只有 pure massive bubble 与 mix bubble+tree 进入外部 reduction 流程。reference bubble 仍先按 reference 顺序由 package \texttt{symmetry} 将 R2 canonical 到 R1，不把它作为另一套 sector metadata 下的独立 DE basis。

% ============================================================
\section{Convention 汇总}
% ============================================================

本节汇总 package 中所有约定设置，包括用途和缺省值。

\subsection{积分表示}

\begin{itemize}
  \item \textbf{Head}：统一使用 $J[\{a_v\}, \{\text{pack}_e\}, \{n_{\text{isp}_j}\}]$。所有 sector（top 和 sub-sector）使用同一 Head。
  \item \textbf{massive 完整线 pack}：$\{b_e, n_{e,1}, n_{e,2}\}$（3 元素）。$b_e$ 为整数指标（分母幂次），$n_{e,a} \in \{0,1\}$ 为 building block 导数阶数。
  \item \textbf{massless 完整线 pack}：$\{b_e,n_{e,1},n_{e,2}\}$（3 元素）。端点关系 $F_{01}=-F_{10}$ 和 $F_{11}=-F_{00}$ 在 canonical 层吸收，但公开 root line 槽不退化。
  \item \textbf{fixed 完整/缩并 pack}：分别为 $\{\texttt{"F"},n_{e,1},n_{e,2}\}$ 与 $\{\texttt{"F"}\}$；字符串只标记没有 $b/bS$ 指标，模长幂属于 sector prefactor。
  \item \textbf{cycle 缩并线 pack}：$\{bS_e\}$（1 元素）。$bS_e$ 为整数指标。
  \item \textbf{$a$ 指标}：正幂次，$(-\tau_v)^{a_v + a0_v}$。
  \item \textbf{$b$ 指标}：分母幂次，$q_e^{-(b_e + b0_e)}$。
\end{itemize}

\subsection{指标零点}

\begin{itemize}
  \item $a0_v$：顶点 $v$ 的 $a$ 零点。h 模式缺省 $2\nu_{\text{ref}}$，H 模式缺省 $0$。
  \item $b0_e$：线 $e$ 的 $b$ 零点。h 模式缺省 $-2\nu_e$，H 模式缺省 $0$。
  \item $bS0_e$：缩并线 $e$ 的 $b$ 零点。h 模式：$bS0_e = b0_e + 2\nu_e$；H 模式：$bS0_e = b0_e$（Wronskian 给出纯 $1/z$，无非整数部分）。
  \item $a0_{\text{merged}}$：合并顶点的 $a$ 零点。h 模式：$a0_{\text{merged}} = a0_u + a0_v - 2\nu_e$；H 模式：$a0_{\text{merged}} = a0_u + a0_v$。
\end{itemize}

\subsection{Building Block 类型}

\begin{itemize}
  \item $\text{bbType}_e$：线 $e$ 的 building block 类型。
  \item \texttt{"h"} $\to \{2\nu+1, 1, -(2\nu+1)\}$：h 函数（归一化 Hankel）。
  \item \texttt{"H"}：裸 Hankel 矩阵 EOM，含 $\nu^2/x^2$ 二次 pole；shrink power 为 $-1$。
  \item \texttt{\{c1,c2\}} 或 line-local \texttt{eomCoefficients}：旧输入兼容；初始化时转换为 $P=c_1/x,Q=c_2,T=I_2,W=-\mathcal Cx^{-c_1}$。
\end{itemize}

\subsection{维度参数}

\begin{itemize}
  \item $d_{\text{PQ}}$：EOM 中的维度参数。缺省 $3$。与维正规化维度独立。
  \item $d$：维正规化维度，出现在动量 IBP 的维度项 $d \cdot J$ 中。缺省 $3$（树图）或 $3-2\epsilon$（圈图）。
\end{itemize}

\subsection{外部能量与 SK 符号}

\begin{itemize}
  \item $P[v]$：顶点 $v$ 的外部能量符号。保持独立到 \texttt{repvar} 阶段。
  \item SK 符号约定：$(+)$ 顶点对应 $e^{-iP_v \tau_v}$，$(-)$ 顶点对应 $e^{+iP_v \tau_v}$。
  \item 时间 IBP 各项符号：vertex power $-a_v$，external energy $-iP_v$，building block derivative $-1$（两端点均取负号）。
\end{itemize}

\subsection{缩并约定}

\begin{itemize}
  \item 缩并条件：h/H 与 masslessFull 均为 $n_{e,1}+n_{e,2}=1$；massless 第一/第二有序端点的 theta-delta 系数分别为 $-2/+2$。
  \item 缩并因子幂次（h 模式）：$-(2\nu_e+1)$。整数 $-1$ 进入指标，非整数 $-2\nu_e$ 进入零点。
  \item 缩并因子幂次（H 模式）：$-1$（纯整数，来自 $W[H_\nu^{(1)}, H_{\nu^*}^{(2)}] \propto 1/z$）。零点无移位。
  \item 缩并因子幂次（无质量）：没有额外 $1/q$、$1/(-\tau)$ 或 $2\nu$ shift；theta-delta contact 把 odd endpoint pack $\{b_e,n_{e,1},n_{e,2}\}$ 变成 $\{bS_e\}$。
  \item $a$ 指标移位（h/H）：$a_{\text{merged}} = a_u + a_v - 1$。
  \item $a$ 指标移位（masslessFull）：$a_{\text{merged}} = a_u + a_v$。
  \item $a$ 零点移位（h 模式）：$a0_{\text{merged}} = a0_u + a0_v - 2\nu_e$。
  \item $a$ 零点移位（H 模式）：$a0_{\text{merged}} = a0_u + a0_v$（无非整数部分）。
  \item $b$ 指标移位（h/H）：$bS_e=b_e+1$。
  \item $b$ 指标移位（masslessFull）：$bS_e=b_e$。
  \item $b$ 零点移位（h 模式）：$bS0_e = b0_e + 2\nu_e$。
  \item $b$ 零点移位（H 模式）：$bS0_e = b0_e$（无非整数部分）。
  \item 常数 prefactor（h 模式）：$\mathcal{C}_e^{(h)} = \frac{4i}{\pi} e^{\pi \text{Im}[\nu_e]}$。
  \item 常数 prefactor（H 模式）：$\mathcal{C}_e^{(H)} = \frac{4i}{\pi} e^{\pi \text{Im}[\nu_e]}$（与 h 模式相同，来自 $W[H_\nu^{(1)}, H_{\nu^*}^{(2)}] = -e^{\pi \text{Im}[\nu]} \frac{4i}{\pi z}$，已数值验证）。
  \item 常数 prefactor（无质量）：$\mathcal{C}_e = 1$（无 prefactor）。
  \item 同一 contact 事件选中 $k$ 条线时：系数另乘 $2^{1-k}$，prefactor 为所选线 $\mathcal C_e$ 的乘积；整个事件仍只有一个 delta。
  \item Sub-sector 中塌缩线的 $\nu$ 消失：$\nu_e \to 0$（参考代码约定）。
\end{itemize}

\subsection{符号选用}

\begin{itemize}
  \item 不同物理对象使用不同符号（如动量 $k_i$ 与指标 $i$ 不混同）。
  \item 当指标不在同一公式中与物理量冲突时，优先使用 $i, j, k$ 作为指标。
  \item 圈动量符号：$q_l$（$l$ 为圈编号）。
  \item 外部能量符号：$P_v$（$v$ 为顶点编号）或 $k_v$。
\end{itemize}

% ============================================================
\section{Loop reduction 取回、微分方程与标度检验}
% ============================================================

018 继续保持 ``package 不运行 reduction'' 的边界，但允许读回用户已经运行完成的 Kira 全套输出。导出端的 \texttt{linearData}、积分编号映射和 convention manifest 必须与读回结果同源；否则 reduction rule 即使语法可读，也不能用于微分方程。

\subsection{用户定义主积分坐标 \texttt{userMI}}

物理积分对象始终只有 \texttt{J}。公开 token \texttt{userMI[i]} 只表示用户在同一 \texttt{J} 线性空间中选定的第 $i$ 个有序坐标，不携带新的 sector 或积分 convention。用户在 \texttt{DSLinear} 之后调用
\begin{verbatim}
prepared = DSUserMI[linearData, expressions, spec]
\end{verbatim}
其中 \texttt{expressions} 的每一项必须是单个 \texttt{J} 或 \texttt{J} 的齐次线性组合；\texttt{spec} 可给出 \texttt{names}、\texttt{activeIndices}、\texttt{derivativeVariables} 与 \texttt{scalingDegrees}。若用户需要改变原始 \texttt{J} 编号，必须先调用 \texttt{DSReorderIntegrals}，再构造 \texttt{userMI}；basis 关系附加后禁止再次重排。

设全部候选组合和其 support 积分分别为 $U_i$ 与 $J_a$，程序精确构造
\begin{equation}
 U_i=\sum_a C_{ia}J_a.
\end{equation}
全体候选行与 \texttt{activeIndices} 子集都必须满行秩。程序从 $C$ 选定一组 pivot 列 $P$，把其余列记为 spectator 集 $S$，并保存
\begin{equation}
 J_P=C_P^{-1}\left(U-C_SJ_S\right).
\end{equation}
因此“可逆”只声称在候选 support 上用 $U$ 替换 pivot $J$、同时保持 spectator $J$ 不变；不把少于全局 \texttt{integralList} 维数的用户 basis 伪装成整个积分空间的方阵变换。manifest 保存系数矩阵、秩、pivot/spectator、双向规则、两向 round-trip residual、原 \texttt{integralList} 顺序 digest，以及 \texttt{userMI[i]} 与 backend \texttt{Tuserweight[i]} 的双向映射。解析一阶导数、最小 target closure、import 和 \texttt{DSDE} 继续复用同一个 active-basis 实现；import 的公开 \texttt{masterTokens} 使用 \texttt{userMI[i]}，后端 token 单独保存在 \texttt{backendMasterTokens}。

\subsection{Kira 内部的实能量变量}

公开物理层仍使用用户初始化的相位能量 $k$。仅在 Kira serializer 内部，对初始化 phase-energy dependency data 中实际进入顶点相位的每个独立原子引入一个实后端符号 $ik$，并取
\begin{equation}
 k=-i\,ik,\qquad ik=i\,k.
 \label{eq:kira-real-energy-map}
\end{equation}
这里的字符 \texttt{ik} 是单个小写 backend symbol，不是 Mathematica 的乘积 \texttt{I*k}。复合顶点能量先依照初始化保存的依赖结构拆成独立原子；同一原子在多个顶点复用时只生成一个 backend symbol，并合并来源顶点和角色。没有进入任何顶点相位的 Gram/根号空间坐标不作此替换。实现不得按符号名称猜测能量；非原子、保留名、大小写折叠重名和与已有 coefficient symbol 的碰撞均令 export fail closed。

普通导数由链式法则给出
\begin{equation}
 {\partial\over\partial k}=i{\partial\over\partial ik},
 \qquad
 {\partial\over\partial ik}=-i{\partial\over\partial k},
 \label{eq:kira-real-energy-jacobian}
\end{equation}
而 Euler 算符严格不变：
\begin{equation}
 k{\partial\over\partial k}=ik{\partial\over\partial ik}.
 \label{eq:kira-real-energy-euler}
\end{equation}
因此 Kira 内部可以全程以实 \texttt{ik} 做 reduction 和 scaling；取回普通 DE 时必须按式~\eqref{eq:kira-real-energy-jacobian} 恢复虚系数。若 $A_k$ 与 $A_{ik}$ 分别表示物理和后端导数矩阵，则
\begin{equation}
 A_{ik}=-iA_k,\qquad A_k=iA_{ik}.
\end{equation}

在 \texttt{postDerivative} 路线中，程序先冻结物理 active basis 的解析一阶导数和最小 target closure，再处理数值规则。物理规则 $k\to r$ 被写成 backend 截面 $ik\to r$，其中 $r\in\mathbb Q$ 且为实数；manifest 同时保存物理求值截面 $k\to-ir$。例如 \texttt{P0->29/13} 生成 \texttt{ip0->29/13} 与 \texttt{P0->-29 I/13}。export 的顺序是物理能量映射、backend 有理数值代入、剩余 Gaussian integral phase gauge；importer 依次恢复一般 coefficient map、backend energy map 和积分 phase gauge。Gaussian gauge 只允许作为残余相位的后备，并保存全部可逆相位数据；不能代替式~\eqref{eq:kira-real-energy-map}。

reference bubble 的数值对照不重新生成 reference IBP 或运行 reference Kira，而是复制并按固定 SHA-256 核验原始 reference 已导出的解析矩阵 \texttt{DEP0.m}、\texttt{DEks.m}、\texttt{DEscaleCheck.m}、\texttt{MIdlogNote.m} 与 \texttt{derivative\_rules\_bubble.m}。设 package 与 reference 的物理能量分别为 $P_{\rm pkg}$、$P_{\rm ref}$，当前 branch 的结构映射为
\begin{equation}
 P_{\rm pkg}=-P_{\rm ref},\qquad
 \frac{\partial}{\partial P_{\rm pkg}}=-\frac{\partial}{\partial P_{\rm ref}}.
\end{equation}
因此 package 点 $P_{\rm pkg}=-29i/13$ 对应 reference 解析矩阵中的 $P_{\rm ref}=29i/13$。解析 stored matrix 先用齐次 degree 做 $N A N^{-1}$，不加入 $N'N^{-1}$；再以 $T=\operatorname{diag}(1^{14},k_s^4,1)$ 恢复 \texttt{MIdlogNote} 第 15--18 项的显式 $k_s$，其中 $A_{k_s}$ 必须加入 $T'T^{-1}$。最后由 $A_{ip0}=-iA_{P0}$ 得到 backend view。19 个 master 定义逐项比例均为 $1$，physical $P_0$、backend $ip0$ 与 physical $k_s$ 三套 $19\times19$ 矩阵均有 $361/361$ 个矩阵元精确相等。

为检验这 19 个 master 不是有限撒点边界留下的假 master，package 侧还必须在不触碰 reference producer 的条件下扩大目标包络。018 验收把两项 $a$ 从 $[-1,4]$ 扩为 $[-2,5]$，两项 $b$ 从 $[-2,5]$ 扩为 $[-3,6]$。fresh 系统含 197232 条生成关系、7850 个 linear 积分、48699 个 Kira 方程和 7871 个导出积分；formal active basis 仍为 19 项。215 个 Kira targets 的精确构成为 19 个 active-basis token 加 196 个 derivative $J$，旧系统的 196 个 derivative 对象在扩张系统中缺失数为零。Kira 2.3 仍选出 backend IDs $1,\ldots,19$，全部 215 targets 约化且 unreduced 为零；回读后的 scaling matrix/source residual 为零，上述三套矩阵再次各通过 $361/361$。

设完全同序的主积分列向量为 $\bm I=(I_1,\ldots,I_N)^T$。对外部独立变量 $x$，公开总导数 \texttt{ds[expr,x]} 同时作用于显式系数和积分：
\begin{equation}
 \frac{d}{dx}\sum_i c_i(x)I_i(x)
 =\sum_i c_i'(x)I_i(x)+\sum_i c_i(x)\frac{dI_i(x)}{dx}.
\end{equation}
把右端用导入的 reduction rules 约回同一顺序，得到
\begin{equation}
 \frac{d\bm I}{dx}=\Omega_x\bm I+\bm s_x.
\end{equation}
如果仍有不在 $\bm I$ 中的积分，或者 $\bm s_x$ 不符合声明的非齐次结构，则该变量的 DE 状态必须是 \texttt{notClosed}。

pure massive bubble 的 top sector 使用文献 2604.14549 Eq.~(51) 的 Euler 标度关系
\begin{equation}
 \left(k_s\partial_{k_s}+P_1\partial_{P_1}+P_2\partial_{P_2}\right)
 I[\{n\},\{a_1,a_2\},\{b_1,b_2\}]
 =\left(d-b_1-b_2-a_1-a_2-2\right)I[\cdots].
 \label{eq:loop-top-scaling}
\end{equation}
residual/tadpole sector 使用该文 Eq.~(64)
\begin{equation}
 \left(k_s\partial_{k_s}+P_0\partial_{P_0}\right)
 R_1[\{n_3,n_4\},\{a\},\{b_1,b_2\}]
 =\left(d-b_1-b_2-a-1\right)R_1[\cdots].
 \label{eq:loop-residual-scaling}
\end{equation}
017 的 scaling check 把这些 Euler operator 逐个作用于 master vector，再经相同 reduction rules 约回。一般写成
\begin{equation}
 \partial_{x_a}\bm I=A_a\bm I+\bm s_a,
 \qquad {\cal E}=\sum_a w_a x_a\partial_{x_a},
 \qquad {\cal E}I_i=\delta_i I_i,
\end{equation}
则同序齐次 master basis 必须满足
\begin{equation}
 \sum_a w_a x_a A_a=\operatorname{diag}(\delta_i),
 \qquad \sum_a w_a x_a\bm s_a=0.
\end{equation}
对一般 loop topology，\texttt{DSScaleCheck} 的 \texttt{LoopTopology} 路线逐个 master 选择其目标 sector，并使用
\begin{equation}
 \delta\!\left[J_s\right]
 =L_{\rm root}d-\sum_{v\in s}A_{v,s}-|V_s|-\sum_e B_{e,s}
  +2\sum_j r_{j,s}+\frac{{\cal E}N_s}{N_s}.
 \label{eq:general-loop-scaling-degree}
\end{equation}
这里 $A_{v,s}=a_{v,s}+a_{0,v,s}$；full line 使用 $B_{e,s}=b_{e,s}+b_{0,e,s}$，shrunk line 使用 $B_{e,s}=b_{S,e,s}+b_{S0,e,s}$；$r_{j,s}$ 是 ISP 分子幂次，因为每个 ISP 是二次齐次的 loop scalar product，所以贡献 $2r_{j,s}$。$L_{\rm root}$ 来自 root topology；contact sector 不重新降圈。最后一项必须从结构化 \texttt{sectorPrefactorData} 物化完整 $N_s$ 后求 Euler 对数导数，不能只看 \texttt{J} 的 pack。若 $N_s$ 或 master 的显式线性组合关于所声明的 Euler 变量不齐次，该路线返回无效 scaling specification，而不是选择一个次数。\texttt{PureMassiveBubble} 仍保留为 reference 专用关系，不负责含 ISP 的一般 topology。

若 $\bm I'=T(x)\bm I$，则
\begin{equation}
 A'_a=(\partial_{x_a}T)T^{-1}+TA_aT^{-1},\qquad
 B'={\cal E}[T]T^{-1}+TBT^{-1},
 \quad B=\sum_a w_a x_aA_a.
\end{equation}
因此依赖动力学变量的 normalization 不能只做共轭；对 $M_i=N_iJ_i$，次数还必须包含 ${\cal E}[N_i]/N_i$。检查对象是上述全符号矩阵和 source 残差；数值点只能定位失败，不能取代式~\eqref{eq:loop-top-scaling}、\eqref{eq:loop-residual-scaling} 及一般 Euler 恒等式的符号检验。

用户用 helper 函数构造 \texttt{symmetryRules} 时，序列化的 topology/checkpoint 保存的是规则表达式，不会自动保存 helper 的 DownValues。因此从 checkpoint 新开 kernel 续跑 \texttt{DSKiraExport} 或 \texttt{DSDE} 前，consumer 必须先加载生成该 topology 的同源 family convention；否则规则头仍存在但条件或右端 helper 无法求值，顶点交换等价的对象会错误地表现成不同 derivative targets。正常单 kernel 全流程没有这一步；只有显式 checkpoint round trip 需要恢复用户定义环境。

reference bubble 还有一个检查侧边界。原始 \texttt{MIdlogNote.m} 的第 15--18 个 dlog 基元显式含 $k_s$，但 \texttt{002 bubble\_de.m} 构造 \texttt{MIdlogSym} 时经 \texttt{reppara2N} 取 $k_s\to1$。因此由 \texttt{DEP0.m/DEks.m} 提取的 stored matrix 即使按齐次次数恢复到一般 $k_s$，仍不是含显式 $k_s$ 的 physical dlog basis。最终对照必须令
\[
 T=\operatorname{diag}(1^{14},k_s^4,1),\qquad
 A^{\rm phys}_{P_0}=T A^{\rm stored}_{P_0}T^{-1},\qquad
 A^{\rm phys}_{k_s}=(\partial_{k_s}T)T^{-1}+T A^{\rm stored}_{k_s}T^{-1}.
\]
相应四个 master 的 Euler degree 各增加一。这个步骤只恢复 reference 源码中原本存在的 $k_s$，不改变 package 的 normalized-$J$ 定义，也不允许从最终差矩阵反解其它 basis adapter。

闭环 example 固定选 pure massive bubble reference 的 $--$ branch，并使用与 reference code 相同的 parity-closed subsystem：top sector 中 $n_1+n_2+b_1$ 与 $n_3+n_4+b_2$ 均为偶数；$R_1$ 中 $b_1$ 与 $n_3+n_4+b_2$ 均为偶数。交换对称性、$R_2\to R_1$ 与 parity 必须在前端统一 canonical 中应用，不能在 Kira importer 或 DE 层临时补入。

% ============================================================
\section{Tree vertex family 的指标表示}
% ============================================================

Tree 与 loop 共用积分 Head，但用 arity 区分：
\begin{align}
 J[\bm a,\bm p,\bm r]&\qquad\text{loop 三槽表示},\\
 J[\{\bm v_1,\ldots,\bm v_E\}]&\qquad\text{tree 单槽表示},
\end{align}
其中
\begin{equation}
 \bm v_e=\{a_e,n_{e1},\ldots,n_{e p_e}\}.
 \label{eq:tree-vertex-pack}
\end{equation}
$a_e$ 是顶点时间幂次的整数偏移；$p_e$ 是该顶点所连、有标记的 massive $h$ 外腿数；$n_{ei}\in\{0,1\}$ 是对应 $h$ 的导数态。massless 外腿进入顶点能量 $k_{0;e}$，但不占 $n_{ei}$ 槽。初始化必须验证每个 pack 的长度为 $1+p_e$。

对单个 $p$-fold vertex family，主积分按 2401.00129 Eq.~(3.33) 排序：
\begin{equation}
 j=1+\sum_{i=1}^{p}n_i2^{p-i}.
 \label{eq:tree-master-order}
\end{equation}
因此最后一个二进制指标变化最快，master list 为 $\{00\cdots0,00\cdots1,\ldots,11\cdots1\}$，共有 $2^p$ 个。任何 dlog 矩阵必须与按式~\eqref{eq:tree-master-order} 生成的 master list 一同保存。

\subsection{Loop 到 tree 的零点与显式系数}

Tree family 不保留 loop 的 $b$ 槽，因此投影必须把完整物理幂次写入显式系数。设目标项属于 sector $s$，原 time seed 的参考积分属于 sector $r$，定义
\begin{align}
 A_s&=\sum_v\left(a_v^{(s)}+a0_v^{(s)}\right),\\
 B_{e,s}&=\begin{cases}
 b_e^{(s)}+b0_e^{(s)},&e\text{ 为 full line},\\
 bS_e^{(s)}+bS0_e^{(s)},&e\text{ 为 shrunk line}.
 \end{cases}
\end{align}
相对参考积分归一化后的投影系数为
\begin{equation}
 \mathcal P_{s\leftarrow r}
 =(-1)^{A_s-A_r}\prod_e k_e^{-\left(B_{e,s}-B_{e,r}\right)}.
 \label{eq:loop-tree-projection-prefactor}
\end{equation}
目标 sector 的 $a0_v^{(s)}$ 同时作为式~\eqref{eq:tree-m10} 的 $\nu_{0;v}$。式~\eqref{eq:loop-tree-projection-prefactor} 必须直接由幂次差构造；对一般符号 $b$ 不使用 \texttt{PowerExpand}。

对 h-mode massive contact，$bS=b+1$、$bS0=b0+2\nu$，同时 $a_{\rm merged}=a_u+a_v-1$、$a0_{\rm merged}=a0_u+a0_v-2\nu$。因此整数和 zero-point 共同产生完整的 $(-k)^{-2\nu-1}$；只写 $k^{-1}$ 会遗漏物理系数。H-mode 采用相同整数 shift 但 zero-point shift 为零，无质量 contact 的两类 shift 均为零。

017 的 sector-tagged \texttt{treeLinearData} 必须为每个未合并 loop 单项分别序列化目标与参考的整数幂、zero-point、完整物理幂及其差值，并保存由这些差值直接重建的式~\eqref{eq:loop-tree-projection-prefactor}。若若干来源投影到相同的 \texttt{\{sectorKey,J[...]\}}，最终系数可以求和，但逐来源 \texttt{contributions} 和 \texttt{physicalPowerAudits} 不得覆盖。对同一 contact 事件同时缩并多条线时，式~\eqref{eq:loop-tree-projection-prefactor} 的乘积逐线保留；例如三条 h 线的显式能量指数分别为 $-1-2\nu_i$，而 merged vertex 的 $\nu_0$ 同时减去 $\sum_i2\nu_i$。

% ============================================================
\section{Tree 的通用迭代关系}
% ============================================================

\textbf{当前实现边界（2026-07-29）}：本节保留原 dSIBP tree 公式的历史 convention 与对照公式；018 中的公式型 \texttt{repIterative}、\texttt{DSTreeDLogDE} 及 massless quotient tree producer继续暂记为禁用，不属于正式公开能力。直接公式实现已经隔离到根目录独立程序包 \texttt{package-MadStree/}：它以原生 time-only context 提供 \texttt{MSInitTree}、\texttt{MSInitTimeGraph} 与专用单顶点入口 \texttt{MSInitVertexFamily}，由局部 $1\times1/2\times2$ building blocks 直接构造主积分、迭代约化和 dlog DE。当前 MadStree 已实现共同-theta odd-subset simultaneous contact、coincident massless/massive quotient、event-reachable sectors、非零 child shift 的公式约化右复合、nested blow-up/theta fixing/normal-crossing 证书，并可调用独立 FlintNDE 后端；它的 time-only 圈图不执行 loop-momentum integration。该隔离不改变本 package 的 topology-driven IBP/serializer 边界，也不把 MadStree 的共享 massless 二态写回旧 \texttt{J} 表示。

定义 $\bm f(a_0)$ 为式~\eqref{eq:tree-master-order} 排序的 $2^p$ 维向量。令 $\Lambda_r^{(i)}$ 表示只在第 $i$ 个二进制因子上作用 Pauli 矩阵 $\sigma_r$。2401.00129 Eq.~(3.37) 给出
\begin{align}
 M_1(\nu_0)
 &=\sum_{i=1}^{p}\left(\nu_i+\frac12\right)\Lambda_3^{(i)}
 +\left(\nu_0-\frac p2-\sum_{i=1}^{p}\nu_i\right)I_{2^p},\\
 M_0
 &=-i\sum_{i=1}^{p}k_i\Lambda_2^{(i)}+ik_0 I_{2^p},
 \label{eq:tree-m10}
\end{align}
并满足
\begin{equation}
 M_1(\nu_0)\bm f(-1)+M_0\bm f(0)=0.
\end{equation}
所以降低一步的矩阵是
\begin{equation}
 A_-(\nu_0)=-M_1(\nu_0)^{-1}M_0.
 \label{eq:tree-aminus}
\end{equation}

为避免一般 $2^p\times2^p$ 矩阵求逆，采用该文 Eq.~(3.43)--(3.50) 的常数变换
\begin{equation}
 T=\frac1{\sqrt2}\begin{pmatrix}1&-i\\-i&1\end{pmatrix},
 \qquad T_p=T^{\otimes p}.
\end{equation}
在变换后基底，
\begin{align}
 \widetilde M_0
 &=-i\sum_{i=1}^{p}k_i\Lambda_3^{(i)}+ik_0I_{2^p},\\
 \widetilde M_1&=T_pM_1T_p^{-1},
\end{align}
其中 $M_1$ 和 $\widetilde M_0$ 都是对角矩阵。升高一步为
\begin{equation}
 A_+(\nu_0-1)
 =-T_p^{-1}\widetilde M_0^{-1}\widetilde M_1T_p
 =-T_p^{-1}\widetilde M_0^{-1}T_pM_1(\nu_0).
 \label{eq:tree-aplus}
\end{equation}
实现应直接构造两个对角逆；当 $M_1$ 或 $\widetilde M_0$ 的对角元为零时返回 singular-locus 信息，不调用大矩阵符号 \texttt{Inverse}。

原始可替换单步规则命名为 \texttt{repIterative0}。函数
\begin{verbatim}
repIterative[expr_, end_:Automatic]
\end{verbatim}
把每个顶点的 $a_e$ 迭代到指定终点；\texttt{Automatic} 表示全零终点。显式终点的长度必须等于顶点数，且每段步数必须是可判定整数。多 sector 情形按 contact-reachable sector DAG 从低到高处理，避免把含 source 的非齐次关系当作齐次 replacement repeated。017 的公开输入始终是三参数 $J$；massive-only family 先按 sector tag 映射到 Private vertex basis，公式内核完成单步关系后再映回同一公开表示。这个要求保证用户显式给出的 \texttt{treeEnergy} 不会被 fixed-line notation \texttt{sEi} 替换。含 massless full line 时 quotient-basis 公式尚未重推，接口返回 \texttt{PendingRederivation}，不使用旧单槽公式猜测结果。

% ============================================================
\section{Tree 的 dlog DE 与 SK branch}
% ============================================================

定义 2401.00129 Eq.~(3.54) 的两个矩阵值对数
\begin{align}
 (\widetilde\Omega_0)_{\bm b\bm a}
 &=\delta_{\bm b\bm a}\left[-i\log\left(k_0+\sum_i(2a_i-1)k_i\right)\right],\\
 (\Omega_{\mathrm{ex}})_{\bm b\bm a}
 &=\delta_{\bm b\bm a}\left[-\sum_i a_i(2\nu_i+1)\log k_i\right].
\end{align}
则 master vector 的通用 dlog connection 是该文 Eq.~(3.55)
\begin{equation}
 d\bm f=(d\Omega)\bm f,
 \qquad
 \Omega=\Omega_{\mathrm{ex}}
 -iT_p^{-1}\widetilde\Omega_0T_pM_1(\nu_0+1).
 \label{eq:tree-dlog}
\end{equation}
013 输出式~\eqref{eq:tree-dlog} 时必须同时给出 letters、矩阵和式~\eqref{eq:tree-master-order} 的 master list。序列化时按顶点顺序逐块输出：每块先列该顶点 massive legs 顺序的能量 letters，再按二进制 master 顺序列 cut letters，最后稳定去重；letter matrices 的键顺序必须与此完全一致，不能依赖展开后表达式的遍历顺序。

\subsection{Naive time-IBP 路线的微分方程}

017 另保留一条不使用式~\eqref{eq:tree-dlog} 或 $A_\pm$ 构造约化规则的路线。对每个 sector 的 $a_v=0$ binary master，把一个活动顶点改成 $a_v=1$，反投影为 loop 代表元并调用已经验证的 \texttt{dtau}。完成共同 theta、simultaneous contact、zero point 和 coincident canonical 后再投影为 sector-tagged tree 方程。所有方程一次联立，固定 $a_v=0$ masters，只解一步升幂对象。这正是 \texttt{DSTreeNaiveIBP} 的定义。

顶点相位能量导数仍可由 loop 代表元的相位项投影。传播子的 \texttt{treeEnergy} 则不能使用完整 loop 动量导数，因为后者作用于积分变量。由式~\eqref{dkh0}--\eqref{dkh1}，以 tree 的 $\tau^a$ convention 写成
\begin{align}
 \partial_k J[a,n=0]&=-J[a+1,n=1],\label{eq:tree-naive-dk0}\\
 \partial_k J[a,n=1]&=J[a+1,n=0]-\frac{2\nu+1}{k}J[a,n=1].\label{eq:tree-naive-dk1}
\end{align}
每个 endpoint leg 分别应用式~\eqref{eq:tree-naive-dk0}--\eqref{eq:tree-naive-dk1}，并乘 $\partial k/\partial x$。这里 $J_s=N_sI_s$ 已隐含 sector prefactor，故
\begin{equation}
 \partial_xJ_s=(\partial_x\log N_s)J_s+N_s\partial_xI_s,
\end{equation}
所以 Wronskian/zero-point prefactor 的导数不可遗漏，也不得在 $J_s$ 外再次乘 $N_s$。\texttt{DSTreeNaiveDE} 最后只使用上述 raw derivative 与 naive 线性规则，输出的逐变量矩阵应与 $\partial_x\Omega$ 在完全相同的 ordered normalized basis 下逐项相等。

每个顶点的 $+/-$ branch 决定其相位能量 $k_{0;e}$ 的符号。h 的 EOM、二进制 basis 和式~\eqref{eq:tree-m10} 的矩阵结构不因 branch 改变。内部线的 SK 类型由两个端点的 branch 唯一决定：同号端点只能给 $G^{++}$ 或 $G^{--}$，异号端点只能给 $G^{+-}$ 或 $G^{-+}$，不存在同一传播子在一次计算中同时具有两种类型的情形。

$G^{+-}/G^{-+}$ 不含 theta，因此 time-IBP 完全按 vertex family 因子化，tree 映射不得读取 $W_T$。$G^{++}/G^{--}$ 的 theta 导数产生 contact remaining term；其 branch 符号必须从当前 loop time-IBP 的 compiled $W_T$ 和 theta offset 读取。对同一顶点对的多条 full lines，先在 loop 表示中应用共同-theta odd-subset contact、simultaneous shift、zero point 和 coincident canonical，再映射到 tree。这样多顶点 DE 按 lower sector 形成 block-triangular source，正是 2401.00129 Eq.~(3.66)--(3.68) 的结构，而不需要在 tree 模块重写传播子的 theta 公式。

设 sector $s$ 的 binary master vector 为 $\bm f_s^{(0)}$。非对角块需要约化求导产生的 $\bm f_s^{(1)}$，所以 contact selector 必须从 $R^{(1)}$ 取得：对 $\bm f_s^{(0)}$ 的每个 binary state，只把当前求导顶点的整数时间指标改成 $a=1$ 后调用 loop time-IBP。$a=0$ seed 给出的是 $R^{(0)}$；其 lower-sector 指标一般为 $a=-1$，先把它迭代到零会引入合并顶点能量，不能替代 Eq.~(3.68) 的 selector。

每个 sector 的 binary master 已按 $\bm J_s=N_s\bm I_s$ 定义；$N_s$ 来自传播子收缩并保存在 sector metadata 中，不再引入第二层 master normalization。于是对角与非对角块都使用同一份 $N_s$：
\begin{align}
 \Omega_{ss}^{\rm global}&=\Omega_s+\log(\mathcal N_s)\,\mathbb I_s,\\
 \Omega_{st}^{\rm global}&=-i\sum_{v\in s}
 \operatorname{Embed}_v\!\left(T_v^{-1}\widetilde\Omega_{0;v}T_v\right)
 C_{st}^{(v)},\qquad t\prec s,
 \label{eq:tree-sector-offdiagonal}
\end{align}
其中 $C_{st}^{(v)}$ 由当前 loop \texttt{dtau} 的 tagged $R^{(1)}$ raw source 乘 $N_s/N_t$ 得到。该写法逐项实现 Eq.~(3.68)，且共同-theta多传播子 source 会自然给直接跨多层 sector 的块。$\log(N_s)$ 作为整体 letter 保存；不把一般复幂拆开，也不使用 \texttt{PowerExpand}。

\texttt{DSTreeDLogDE[context]} 返回同序的 \texttt{masters}、\texttt{omega}、\texttt{letters/letterMatrices}，并附带 \texttt{sectorNormalizations}、\texttt{normalizationAudits}、\texttt{contactMaps} 与 \texttt{omegaBlocks}。缺省 \texttt{standard} 路线信任同一 producer 的统一公开表示 stamp，不再全对象扫描 legacy/private token，也不构造 \texttt{dlogResidual}；显式 \texttt{AuditLevel->"full"} 才运行公开表示、\texttt{LinearSolve} residual 与 dlog 严格重构证书。矩阵按 contact DAG 排序；从 lower sector 回到 parent 的块及同层 sibling 块必须严格为零。

% ============================================================
\section{013 pure time-IBP 交叉验证定义}
% ============================================================

013 的新增 benchmark 不重复已有 momentum-IBP、独立变量总导数或共同-theta专项。它固定使用两个只含 time generator 的 family：
\begin{enumerate}
  \item 两顶点同号 $\{+,+\}$，由一条 massive full line 相连。该线是 $G^{++}$，必须出现由 compiled $W_T$ 决定的 contact source。
  \item 三顶点 chain $\{+,+,-\}$，两条 massive full lines 分别连接 $(1,2)$ 与 $(2,3)$。第一条为 $G^{++}$，第二条为 $G^{+-}$；只有第一条允许产生 contact source。
\end{enumerate}
每个顶点可连 massless 外腿，但其总能量 $K_e$ 必须作为与内部 massive 线能量 $k_{12},k_{23}$ 独立的符号输入。这样 time-IBP 的顶点相位项与内部 building-block 导数项不会因误用同一个变量而发生伪抵消。

对每个 case，验证顺序固定为：先由式~\eqref{eq:tree-m10} 和 loop time-IBP/contact 公式手推 general seed；再调用 package 的 \texttt{dtau}；随后把 seed 投影为 tree 单步关系，与式~\eqref{eq:tree-aminus}、\eqref{eq:tree-aplus} 生成的 \texttt{repIterative0} 比较。最后选择避开 $M_1$、$\widetilde M_0$ 奇异面的有理参数点，把 lower-sector masters 和 top masters 赋予固定有理数，分别用 seed 解与 \texttt{repIterative} 约化到同一终点。两路结果之差必须严格为零。

即使 case 用 loop topology 描述，013 benchmark 也只调用 \texttt{dtau}，不生成或消费 \texttt{dqq}/\texttt{dqk}。这把检验范围限制在本版本新增的 pure time-IBP/tree 功能。

\appendix

% ============================================================
\section{共同-theta 方案的完整推导}
\label{app:commonThetaDerivation}
% ============================================================

设有限条 full lines 共享同一变量 $\Delta$，并满足式~\eqref{bundleSingleLine}。在 $\Delta>0$ 的开区域，每条线都取 $A_e$；在 $\Delta<0$ 的开区域，每条线都取 $B_e$。任何同时含 $\theta(\Delta)$ 与 $\theta(-\Delta)$ 的交叉项只可能位于单点 $\Delta=0$，作为局部可积函数按 almost-everywhere 意义为零。因此乘积作为 distribution 与式~\eqref{bundleCommonThetaProduct} 相同。

对其求分布导数，regular 部分来自 $A_e,B_e$ 的普通导数；奇异部分为
\begin{equation}
 \partial_\Delta\theta(\Delta)\prod_eA_e
 +\partial_\Delta\theta(-\Delta)\prod_eB_e
 =\delta(\Delta)\left(\prod_eA_e-\prod_eB_e\right).
 \label{appCommonBoundary}
\end{equation}
所以整个 bundle 只有一个 delta。该结论先于且独立于任何 $\theta(0)$ 点值约定。

由式~\eqref{bundleJDDefinition}，
\begin{equation}
 A_e=\mathsf J_e+\frac{D_e}{2},
 \qquad
 B_e=\mathsf J_e-\frac{D_e}{2}.
\end{equation}
展开两个乘积时，选择子集 $S$ 表示哪些因子取 $D_e$。两个乘积中同一子集的系数之差为
\begin{equation}
 \left(\frac12\right)^{|S|}
 -\left(-\frac12\right)^{|S|}
 =
 \begin{cases}
 0,&|S|\text{ 为偶数},\\
 2^{1-|S|},&|S|\text{ 为奇数}.
 \end{cases}
\end{equation}
这就得到式~\eqref{commonThetaOddSubset}。偶数 contact 的消失是左右区域交换下的奇偶性，不是额外 sector 规则。

若采用对称点值 $\theta(0)=1/2$，单线等时值是 $\mathsf J_e$。但式~\eqref{appCommonBoundary} 不能通过把每个未求导 theta 独立替换为 $1/2$ 得到；这种替换会丢失三个及以上传播子中的高阶奇数 $D$ 项。正确的指标映射必须继续使用正文式~\eqref{massiveDToWT}--\eqref{bundleFinalJIndexForm}。

\subsection{三线 worked examples}

三条线共享同一 $\Delta$ 时，直接代入 $A_e=\mathsf J_e+D_e/2$、$B_e=\mathsf J_e-D_e/2$ 得
\begin{align}
 A_1A_2A_3-B_1B_2B_3
 ={}&D_1\mathsf J_2\mathsf J_3
 +\mathsf J_1D_2\mathsf J_3
 +\mathsf J_1\mathsf J_2D_3
 +\frac14D_1D_2D_3.
 \label{threeLineDirectExpansion}
\end{align}
前三项分别选择一条线 contact，最后一项同时选择三条线；四项都乘同一个 $\delta(\Delta)$。

先取三条同向 masslessFull 线，均取 canonical odd endpoint state $\{n_{e,1},n_{e,2}\}=\{1,0\}$，并从它们的第一有序端点求导。由 $\mathsf J_e=0,D_e=-2$，三个 single 项均因含未选中线的 $\mathsf J=0$ 而消失，只剩
\begin{equation}
 \frac14D_1D_2D_3=\frac14(-2)^3=-2.
\end{equation}
因此
\begin{equation}
 J\!\left[\{a_u,a_v\},
   \{\{b_1,1,0\},\{b_2,1,0\},\{b_3,1,0\}\},\mathrm{ispList}\right]
 \longrightarrow
 -2J\!\left[\{a_u+a_v\},
   \{\{b_1\},\{b_2\},\{b_3\}\},\mathrm{ispList}\right].
 \label{threeMasslessIndexExample}
\end{equation}
没有 $a/b$ 整数或 zero-point shift；三个 line packs 同时变为 shrunk，但顶点只合并一次。

再取 line 1 为 h-mode massiveFull，line 2、3 为同向 masslessFull。令 line 1 在 $++$ case 的第一端点状态为 $\{n_{1,1},n_{1,2}\}=\{1,0\}$，则
\begin{equation}
 -W_{T,1}=\mathcal C_1x_1^{-2\nu_1-1},
 \qquad D_1\longmapsto\mathcal C_1x_1^{-2\nu_1-1}.
\end{equation}
若两条 massless 线仍取 canonical odd states $\{n_{2,1},n_{2,2}\}=\{n_{3,1},n_{3,2}\}=\{1,0\}$，则 $\mathsf J_2=\mathsf J_3=0$、$D_2=D_3=-2$。式~\eqref{threeLineDirectExpansion} 仍只剩 triple 项，其系数为
\begin{equation}
 \frac14\mathcal C_1(-2)(-2)=\mathcal C_1.
\end{equation}
最终指标与零点是
\begin{align}
 &J\!\left[\{a_u,a_v\},
  \{\{b_1,1,0\},\{b_2,1,0\},\{b_3,1,0\}\},\mathrm{ispList}\right]
 \nonumber\\
 &\qquad\longrightarrow
 \mathcal C_1
 J\!\left[\{a_u+a_v-1\},
  \{\{b_1+1\},\{b_2\},\{b_3\}\},\mathrm{ispList}\right],
 \label{mixedThreeLineIndexExample}
\end{align}
并配套使用
\begin{equation}
 a0_{uv}'=a0_u+a0_v-2\nu_1,
 \qquad bS0_1=b0_1+2\nu_1,
 \qquad bS0_2=b0_2,
 \qquad bS0_3=b0_3.
\end{equation}
这展示了 $W_T$ 的整数 $1$ 与 zero-point $2\nu_1$ 必须同时进入最终 $J$；两条 massless contact 不附加幂次 shift。

% ============================================================
\section{统一 Gaussian 逐线方案及其与共同-theta的等价性}
\label{app:gaussianDerivation}
% ============================================================

取 Gaussian approximate identity
\begin{equation}
 \rho_\epsilon(x)=\frac{e^{-x^2/\epsilon^2}}{\sqrt\pi\,\epsilon},
 \qquad
 H_\epsilon(x)=\int_{-\infty}^{x}\rho_\epsilon(y)\,dy.
 \label{gaussianMollifier}
\end{equation}
则 $H_\epsilon'=\rho_\epsilon$、$H_\epsilon(0)=1/2$，且 $H_\epsilon(\pm\infty)=1,0$。同一 bundle 的所有线必须使用同一个 $H_\epsilon$：
\begin{equation}
 G_{e,\epsilon}=B_e+H_\epsilon D_e.
 \label{regularizedLine}
\end{equation}

先证明基本分布极限。对任意测试函数 $\varphi$ 和连续函数 $F$，令 $h=H_\epsilon(x)$，则 $dh=\rho_\epsilon(x)dx$。过渡层在 $\epsilon\to0$ 时收缩到 $x=0$，所以
\begin{align}
 \lim_{\epsilon\to0}
 \int dx\,\varphi(x)\rho_\epsilon(x)F(H_\epsilon(x))
 &=\varphi(0)\int_0^1F(h)\,dh,\\
 \rho_\epsilon F(H_\epsilon)
 &\longrightarrow
 \delta(x)\int_0^1F(h)\,dh.
 \label{rhoFDistributionLimit}
\end{align}
特别地，
\begin{equation}
 \rho_\epsilon H_\epsilon^m\longrightarrow\frac{1}{m+1}\delta,
 \qquad
 \rho_\epsilon H_\epsilon^r(1-H_\epsilon)^s
 \longrightarrow {\rm B}(r+1,s+1)\delta.
 \label{rhoPowerLimits}
\end{equation}
因此不能把 $H_\epsilon(0)^m=2^{-m}$ 直接乘到 delta 上；delta 正则层扫描的是整个 $0\le h\le1$，不是只取中心点。

对 $P_\epsilon=\prod_eG_{e,\epsilon}$ 求导，并把 $H_\epsilon'$ 归因到第 $i$ 条线，其奇异部分为
\begin{equation}
 {\cal B}_{i,\epsilon}
 =\rho_\epsilon D_i
 \prod_{j\ne i}(B_j+H_\epsilon D_j).
\end{equation}
由式~\eqref{rhoFDistributionLimit}，
\begin{equation}
 {\cal B}_i
 =\delta(x)D_i\int_0^1dh
 \prod_{j\ne i}(B_j+hD_j).
 \label{lineResolvedBoundary}
\end{equation}
这一定义保留了逐传播子 theta，但单条 $i$ 的分配依赖所选共同 regulator；物理上不依赖 regulator 的是所有 $i$ 的总和。

令 $t=h-1/2$，则 $B_j+hD_j=\mathsf J_j+tD_j$。对固定 $i$ 展开得
\begin{equation}
 {\cal B}_i
 =\delta D_i
 \sum_{\substack{T\subseteq\mathcal B\setminus\{i\}\\|T|\ {\rm even}}}
 \frac{1}{2^{|T|}(|T|+1)}
 \prod_{j\in T}D_j
 \prod_{j\notin T\cup\{i\}}\mathsf J_j,
 \label{lineResolvedMomentExpansion}
\end{equation}
因为
\begin{equation}
 \int_{-1/2}^{1/2}t^m dt
 =0\quad(m\text{ 奇}),
 \qquad
 \int_{-1/2}^{1/2}t^m dt
 =\frac{1}{2^m(m+1)}\quad(m\text{ 偶}).
\end{equation}
所以每个逐线项都含奇数个 $D$。

例如三条线时，
\begin{align}
 {\cal B}_1&=\delta\left(D_1\mathsf J_2\mathsf J_3+\frac1{12}D_1D_2D_3\right),\\
 {\cal B}_2&=\delta\left(\mathsf J_1D_2\mathsf J_3+\frac1{12}D_1D_2D_3\right),\\
 {\cal B}_3&=\delta\left(\mathsf J_1\mathsf J_2D_3+\frac1{12}D_1D_2D_3\right).
\end{align}
其中 $1/12=\int_0^1(h-1/2)^2dh$。求和后 triple coefficient 是 $3/12=1/4$。

一般地，固定一个含 $k$ 条线的奇数子集 $S$。在某个 $i\in S$ 的逐线贡献中，其余 $k-1$ 个 $D$ 来自中心矩，系数为
\begin{equation}
 \frac{1}{2^{k-1}k}.
\end{equation}
共有 $k$ 个 $i\in S$ 可以承担 $H_\epsilon'$，故总系数为
\begin{equation}
 k\frac{1}{2^{k-1}k}=2^{1-k},
\end{equation}
与式~\eqref{commonThetaOddSubset} 完全相同。也可以直接观察
\begin{align}
 \sum_iD_i\prod_{j\ne i}(B_j+hD_j)
 &=\frac{d}{dh}\prod_j(B_j+hD_j),\\
 \sum_i{\cal B}_i
 &=\delta\int_0^1dh\frac{d}{dh}\prod_j(B_j+hD_j)
 =\delta\left(\prod_jA_j-\prod_jB_j\right).
\end{align}

因此两种方案对总 boundary 完全等价，前提是同一 bundle 使用同一个 regulator，并保留全部逐线贡献及其中心矩项。package 可以任选其一作为 canonical；012 选择共同-theta方案，因为它无 regulator 与逐线归因依赖，直接给出式~\eqref{commonThetaOddSubset}。不同线使用不同 regulator、只取 $H(0)=1/2$、丢弃高阶奇数项，或把 simultaneous contact 写成 $\delta^k$，都不属于上述等价类。

\section{图论、fixed line 与 pure-time 表示}

\subsection{结构圈数和 root loop space}

令内部传播子构成无向多重图 $G=(V,E)$，连通分量数为 $C$。结构圈数是第一 Betti 数
\begin{equation}
 L=E-V+C.
 \label{graphBettiLoopCount}
\end{equation}
平行边分别计入 $E$，自环各增加一圈，外腿不属于内部边。删除边 $e$ 后若连通分量增加，则 $e$ 是 bridge；否则是 cycle line。式~\eqref{graphBettiLoopCount} 只给拓扑圈数，不能证明用户 routing 正确；还必须检查 loop-coefficient matrix 的秩及其 incidence-cycle support。

三个最小图可直接区分端点 convention：两顶点间一条普通边有 $(E,V,C)=(1,2,1)$，故 $L=0$ 且该边是 bridge；tadpole 必须写成自环 \texttt{endpoints->\{v,v\}}，有 $(1,1,1)$，故 $L=1$、incidence 列为零且该边是 cycle；两顶点三条平行边有 $(3,2,1)$，故 $L=2$ 且三条边都不是 bridge。初始化后的 topology 顶层同时保存 \texttt{graphLoopCount}、\texttt{cycleLineIndices}、\texttt{bridgeLineIndices} 与 \texttt{selfLoopLineIndices}，完整审计仍保存在 \texttt{graphTopologyAudit}。

所有 contact-reachable sector 继承 root topology 的 $L$、loop basis、cycle/bridge 分类和 line-power schema。sector 只改变 full/shrunk 状态、端点代表、zero-point 与 symmetry。特别地，即使一个 contact sector 中所有 cycle lines 都已 shrink，也不得重新按剩余图计算较小的 $L$；否则 triple contact 会丢失原来的 loop-scalar-product 空间和 $b_S$ 指标。

\subsection{cycle pack、fixed pack 与显式幂系数}

full loop 模式下统一积分仍写成
\begin{equation}
 J[\{a_v\},\{P_e\},\{n_{\rm isp}\}].
\end{equation}
但逐线 pack 由 root schema 决定：
\begin{equation}
  P_e=
 \begin{cases}
  \{b_e,n_{e1},n_{e2}\},&e\text{ 为 massive/massless full cycle line},\\
  \{b_e,0,0\},&e\text{ 为 massless cross cycle line},\\
  \{b_{S,e}\},&e\text{ 为 shrunk cycle line},\\
  \{\texttt{"F"},n_{e1},n_{e2}\},&e\text{ 为 massive/massless full fixed line},\\
  \{\texttt{"F"},0,0\},&e\text{ 为 massless cross fixed line},\\
  \{\texttt{"F"}\},&e\text{ 为 shrunk fixed line}.
 \end{cases}
 \label{cycleFixedPackSchema}
\end{equation}
其中 \texttt{"F"} 只标记 fixed schema，不是幂指标。pack pattern 与 root line 顺序共同唯一编码 shrink set；物理幂必须读取 sector metadata。

设某项要求物理分母幂 $B_e\to B_e+\Delta_e$。统一 line-power 操作为
\begin{equation}
 \operatorname{shiftPower}_e(\Delta_e):
 \begin{cases}
  cJ[\ldots,b_e,\ldots]\mapsto cJ[\ldots,b_e+\Delta_e,\ldots],&e\text{ cycle},\\
  cJ[\ldots,P_e,\ldots]\mapsto c\,r_e^{-\Delta_e}J[\ldots,P_e,\ldots],&e\text{ fixed},
 \end{cases}
 \label{shiftLinePowerBranches}
\end{equation}
其中 $r_e=|Q_e|$ 是该 fixed line 已声明的模长。zero-point $B_{0,e}$ 本来就包含在显式因子 $r_e^{-B_{0,e}}$ 中；EOM、time derivative、外模长导数及 Wronskian/contact 引入的每个额外幂都按式~\eqref{shiftLinePowerBranches} 乘到显式系数，不能伪造一个 $b_e$ 指标。

\subsection{momentum IBP 与 time IBP 的作用域}

momentum IBP 的链式法则只对 cycle line 定义 $z_e=\xi_e^2$：
\begin{equation}
 v\cdot\partial_{q_l}F
 =\sum_{e\in E_{\rm cycle}}c_{el}{v\cdot Q_e\over\xi_e}
 {\partial F\over\partial\xi_e}
 +\sum_j(v\cdot\partial_{q_l}\rho_j){\partial F\over\partial\rho_j}.
 \label{cycleOnlyMomentumIBP}
\end{equation}
bridge/fixed line 不含独立 loop flow，所以不进入 $z$ 坐标、复合 momentum generator、$b$ shift 或 momentum building-block derivative。这并不意味着它从积分中消失：其 $r_e^{-B_e}$ 和 building block 仍是显式动力学依赖，按外变量链式法则求导。

time IBP 的作用域不同。对顶点 $v$ 做 $\tau_v$ total derivative 时，所有 incident active lines 的 endpoint building block 都依赖 $\tau_v$，故
\begin{equation}
 \partial_{\tau_v}\prod_{e\ni v}{\cal B}_e
 =\sum_{e\ni v}(\partial_{\tau_v}{\cal B}_e)
  \prod_{f\ni v,f\ne e}{\cal B}_f.
 \label{timeIBPAllIncidentLines}
\end{equation}
这个和式必须包含 cycle 与 fixed lines。fixed line 的 $a/n$ shift 与式~\eqref{shiftLinePowerBranches} 的显式模长因子同时保留；共同-theta bundle、$W_T$ shrink 和 simultaneous contact 也在这一层处理。

\subsection{direct pure-time 模式和能力门禁}

树图缺省、圈图可选的 \texttt{ibpMode->"timeOnly"} 不构造残缺的 loop family。所有 active lines 统一标为 \texttt{fixedCoefficient}。从 020 起，time-only 的唯一公开对象为
\begin{equation}
 J[\texttt{sectorKey},\{a_C\}_C,\{n_r\}_r].
 \label{directPureTimeJ}
\end{equation}
第一槽是式~\eqref{v019TimeOnlySectorBitKey} 的定长字符串；第二槽按当前 sector 的 compact vertex-component 顺序保存时间幂指标；第三槽是 $n_i$ 类离散函数态，不是第二套 sector key。状态槽按 root line 顺序冻结：每条未收缩 massive line 依次贡献两个端点态，每条未收缩 massless quotient line 整边只贡献一个共享态，纯指数、massless cross 与已收缩 line 不贡献槽，也不保留占位。内部 producer 仍可短暂使用 root-ordered line pack，但该对象只存在于 Private 边界，并由 sector metadata 与状态槽 registry 精确双向转换。全部 active line 模长纳入 \texttt{independentExternalMomenta} 审计，且没有 momentum generator、$z$/ISP 闭合或 $b/b_S$ 指标。\texttt{DSSeeds} 完整枚举允许态并立即应用 EOM/quotient/canonical。论文 vertex basis 只作为 Private massive-only 适配器服务于 \texttt{DSTreeSeeds}、\texttt{repIterative}、naive time-IBP/DE 与直接 dlog DE；结果必须映回式~\eqref{directPureTimeJ} 并使用同序 master 列表比较。含 massless full line 的公式型 tree 接口在 quotient basis 重新推导前返回 \texttt{PendingRederivation}。

用户列表和动力学规则的综合状态决定 context capabilities：exact 开放初始化、time IBP、full 模式下的 momentum IBP、求导及唯一反变换；undercomplete 返回 missing/null-space 后拒绝初始化；overcomplete 允许 symbolic seed，但关闭 \texttt{ds/DSDE} 与唯一反变换。\texttt{DSLinear} 还校验 seed 与 context 的 \texttt{inputHash}，serializer 只消费通过该门禁的 backend-neutral \texttt{linearData}。

缺省 notation 为
\begin{equation}
 \texttt{ssij}=\sqrt{\operatorname{sp}(k_i,k_j)},\qquad
 \texttt{sEi}=\sqrt{\operatorname{sp}(p_i,p_i)},
\end{equation}
其中 $k_i$ 与 $p_i$ 分别按两个显式用户列表排序。\texttt{DSRedefineParameters} 只接受覆盖完整基础坐标的规则，并通过 Jacobian 链式法则复用平方 Gram 原子导数；不另写第二套动量微分核心。

\section{017 固定三槽、sector prefactor 与 parity}

本节记录 017 已实施并由模块 smoke 验证的 convention。公开 API、seed、\texttt{linearData}、求导与 tree 结果一律使用三参数表示；不存在可由用户调用的并行 vertex-pack 积分 Head。

\subsection{每条 line 的固定三槽}

017 的公开对象始终为 $J[\texttt{aList},\texttt{linePacks},\texttt{ispList}]$。每条 full root line 固定为三槽；shrunk line 已无端点态，允许退为单槽，但它在 root-ordered \texttt{linePacks} 中的位置永不删除：
\begin{equation}
 \begin{array}{c|cc}
 &\text{massive}&\text{massless}\\ \hline
 \text{cycle full}&\{b,n_1,n_2\}&\{b,n_1,n_2\}\\
 \text{fixed full}&\{\texttt{"F"},n_1,n_2\}&\{\texttt{"F"},n_1,n_2\}\\
 \text{cycle shrunk}&\{b_S\}&\{b_S\}\\
 \text{fixed shrunk}&\{\texttt{"F"}\}&\{\texttt{"F"}\}
 \end{array}
 \label{v017FixedThreeSlotSchema}
\end{equation}
短字符串 \texttt{"F"} 表示 fixed line 没有圈幂次指标。省略首项写成 \texttt{\{,n1,n2\}} 会被 Mathematica 解释为 \texttt{Null} 并产生 \texttt{Syntax::com}；\texttt{\_} 是 pattern \texttt{Blank[]}；\texttt{Nothing} 会直接删掉列表元素，三者都不能作为正式数据。massless 仍保留两个有序端点态；其额外关系属于 non-IBP relation 层，不再通过缩短 pack 表示。

完整的 root-ordered full/shrunk pack pattern 是 shrink set 的隐式、可逆编码。例如第 2、4 条 root line 收缩时，第二和第四位置分别改成相应的单槽 shrunk pack，由位置直接重建集合 $\{2,4\}$；fixed shrunk line 同样必须保留 \texttt{\{"F"\}} 槽。canonical \texttt{sectorKey} 只能由该 pattern 派生，再由 seed、metadata 与 \texttt{linearData} 显式携带作交叉核验。不得删除收缩槽位、压缩 \texttt{linePacks}，也不得用 compact \texttt{aList} 充当 sector 身份。

019 对 \texttt{ibpMode->"timeOnly"} 进一步固定显式 key 的写法。若 root propagator 按输入顺序为
$E=(e_1,\ldots,e_{N_E})$，定义
\begin{equation}
 \operatorname{sectorKey}(s)=q_1q_2\cdots q_{N_E},\qquad
 q_j=\begin{cases}
 0,&e_j\text{ 已收缩},\\
 1,&e_j\text{ 未收缩}.
 \end{cases}
 \label{v019TimeOnlySectorBitKey}
\end{equation}
不能发生 contact shrink 的传播子恒取 $q_j=1$，top sector 是全 $1$ 字符串。该对象的类型是字符串而不是二进制整数，故前导零必须保留；例如四条 root propagator 中收缩第 $1,2,4$ 条时 key 为 \texttt{"0010"}。metadata 同源保存 \texttt{rootLineOrder}、\texttt{width} 和位语义。所有 top/key predicate 必须同时读取 \texttt{ibpMode} 与 root 位宽；在 \texttt{timeOnly} 输入或生成批次中出现 \texttt{"top"}、\texttt{"contracted:..."} 或收缩线名拼接 key 时，初始化直接失败，不作 legacy alias。full-loop 工作流不属于本次接口变化，继续使用既有 key，以避免改动 Kira/reduction artifact identity。

\subsection{020 time-only 状态槽 registry}

对每个 reachable sector $s$，初始化冻结
\begin{equation}
 \mathcal R_s=(r_1,\ldots,r_{N_s}),\qquad
 r_j=(\text{root line position},\text{endpoint/shared slot}).
 \label{v020TimeOnlyStateRegistry}
\end{equation}
公开第三参数 \texttt{stateBits} 的第 $j$ 项就是 $r_j$ 对应的离散指标 $n_i$。它不是第二套 sector key，也不要求按顶点重新分组：massive line 的两个端点各占一项；massless quotient 的两个端点由额外关系缩成整边共享的一项。收缩一条 line 时，它的全部状态槽从 child sector 的 registry 删除，所以 child 不携带零占位。\texttt{sectorKey} 唯一决定采用哪份 $\mathcal R_s$，而 \texttt{stateBits} 只决定该 sector 内的函数态。逆转换把共享 massless 位写入选定 canonical 端点槽，并把被关系消去的另一槽固定为零；因此往返是对 quotient canonical representative 的精确往返。

\texttt{DSSeeds/DSGenerateIBP/DSLinear/dtau/ds}、tree 公式接口、master list 以及表示辅助函数都使用同一转换边界。旧 time-only \texttt{J[aList,linePacks,\{\}]} 在 020 公开入口直接拒绝；full 模式的 \texttt{J[aList,linePacks,ispList]} 不经过该转换。\texttt{rep2innerform/rep2outform} 只作用于显式动力学系数并保护公开 \texttt{J}；\texttt{rep2Integrand} 与 \texttt{symmetry} 在需要 line metadata 时临时转入 Private producer 表示，完成后再回到式~\eqref{directPureTimeJ}。临时逆转换得到的 root-ordered pack 只服务旧 producer；sector normalization 必须继续按公开 \texttt{sectorKey} 读取初始化冻结的 metadata，不能从 pack shape 重新猜测 child zero point 或模长幂。

\subsection{结构化 sector prefactor 与求导}

fixed line 的模长幂不进入式~\eqref{v017FixedThreeSlotSchema} 的第一槽，而属于每个 sector 的 normalization。初始化首先按用户声明的 \texttt{independentExternalMomenta} 顺序冻结稳定编号
\begin{equation}
 \{kE_i\},\qquad \boldsymbol\beta_s=(\beta_{s,1},\ldots,\beta_{s,m}),
 \qquad \mathcal K=(K_1,\ldots,K_m),
 \label{v017SectorPrefactorLists}
\end{equation}
其中 metadata 以 \texttt{kEpower[$\beta_{s,1}$,$\ldots$,$\beta_{s,m}$]} 保存稳定幂向量，以 \texttt{kEParameterExpressions} 保存当前参数表达式 $\mathcal K$。每个 $K_i$ 可以是复合表达式；例如参数重定义可给出 $\mathcal K=(u+v,u)$。无法直接归入 kE 槽的 fixed 模长另以结构化 residual power part 保存。每条已收缩 massive line 的 Wronskian normalization 以 line-indexed contraction part 保存，simultaneous contact 中同一条线只累计一次。metadata 不保存已经乘开的乘积；唯一 materializer 在需要时构造
\begin{equation}
 N_s=c_s\prod_{i=1}^{m}K_i^{\beta_{s,i}}\prod_r R_r^{\gamma_{s,r}},
 \qquad J_s=N_s I_s.
 \label{v017SectorNormalization}
\end{equation}
\texttt{DSRedefineParameters} 只更新式~\eqref{v017SectorPrefactorLists} 的 $\mathcal K$ 与 residual current-expression mapping，不改 kE 编号或 \texttt{kEpower}。求导、contact 和反变换在消费点才物化 $N_s$。020 time-only consumer 总是从 \texttt{sectorKey} 对应的冻结 metadata 取得这份数据；例如 massive line 已收缩的 child sector 若 $N_s\propto s_E^{2\mu}$，则 \texttt{ds[J[s,...],sE]} 必须保留 $(2\mu/s_E)J_s$。由此
\begin{equation}
 \partial_xJ_s=(\partial_x\log N_s)J_s
 +\sum_t c^{\rm raw}_{st}{N_s\over N_t}J_t,
 \qquad
 \partial_xI_s=\sum_t c_{st}I_t.
 \label{v017NormalizedDerivative}
\end{equation}
所以 prefactor 虽然逻辑上属于 $J_s$，其导数和跨 sector 比值仍必须由 \texttt{ds/DSDE} 调取。\texttt{rep2Integrand} 则先构造不含 sector prefactor 的裸 integrand，再把物化后的完整 $N_s$ 乘回；不得依赖 line renderer 偶然重复某一 fixed zero-point 幂。atomic massless lower sector 若 $N_t$ 已含 $s_E^{-\beta}$，contact 项不得再次外乘同一因子。

018 的 E018-15 修复已把 metadata、contact、\texttt{ds/DSDE} 与 \texttt{rep2Integrand} 接到同一 normalized-$J$ 数据流。局部裸积分展开只用于开发定位；最终统一正确性仍由 fresh reduction 后的完整 scaling relation 检查，不增加新的用户阻断条件。

\subsection{fixed-line shrink normalization convention}

fixed/non-loop line 没有 $b/b_S$ 指标，其 contact 模长幂只由 source/target sector 已选择的 zero point 分配，不设置额外整数吸收开关。若 compiled shrink factor 为
\begin{equation}
 -W_{T,e}=\mathcal C_e r_e^{-s_e-z_e}(-\tau)^{-s_e-z_e},
 \qquad s_e\in\mathbb Z,
\end{equation}
且 source/target sector 中该线的 prefactor 分母指数分别为 $B_{s,e}$、$B_{t,e}$。target 的 $N_t$ 同时吸收该线选定的 Wronskian normalization $g_e$；因此裸系数 $c_e r_e^{-s_e-z_e}$ 统一变为
\begin{equation}
 c^{\rm norm}_{s\to t}
 ={c_e\over g_e}r_e^{-\left[s_e+z_e-(B_{t,e}-B_{s,e})\right]}.
 \label{v017FixedShrinkCoefficient}
\end{equation}
共同-theta 系数、端点符号及 $2^{1-|S|}$ 不包含在这里，按原 contact 公式另乘。

缺省 h/H preset 取 $g_e=c_e=\mathcal C_e$，且 sector zero-point 传播为 $B_{t,e}-B_{s,e}=z_e$，所以式~\eqref{v017FixedShrinkCoefficient} 化为 $r_e^{-s_e}$：Wronskian 常数由 target $J_t$ 吸收，zero point 吸收 $z_e$，只留下整数幂。一般多项式 Wronskian 以稳定的首个非零 compiled coefficient 作为 $g_e$，其它 monomial 保留可复核的 $c_e/g_e$。用户若合法改写 source 或 target zero point，式~\eqref{v017FixedShrinkCoefficient} 自动改变 reduction/DE coefficient。package 不提供 \texttt{fixedShrinkConvention}、\texttt{fullContactPower} 或其它平行吸收状态。

zero-point convention 不改变时间幂映射：始终为 $a_t=a_u+a_v-s_e$、$a_{0,t}=a_{0,u}+a_{0,v}-z_e$。现行 preset 配合缺省 $B_t-B_s=z$ 的结果为
\begin{center}
\small
\begin{tabular}{c|c|c|c}
 类型 & $(s_e,z_e)$ & target prefactor exponent & normalized coefficient 的模长部分\\ \hline
 h & $(1,2\nu_e)$ & $B_{s,e}+2\nu_e$ & $r_e^{-1}$\\
 H & $(1,0)$ & $B_{s,e}$ & $r_e^{-1}$\\
 massless & $(0,0)$ & $B_{s,e}$ & $1$
\end{tabular}
\end{center}
h/H 的 $\mathcal C_e=(4i/\pi)e^{\pi\operatorname{Im}\nu_e}$ 进入 target sector 的 contraction part；massless 没有 Wronskian prefactor，端点 contact 的 $-2/+2$ 仍单独保留。初始化 metadata 对每个 fixed transition 保存 $B_s$、$B_t$、zero-point override 来源、compiled $(s,z)$ 与剩余 exponent $s+z-(B_t-B_s)$，从而可由裸系数和 $N_s/N_t$ 唯一重建 normalized 系数。

对 cycle line，现行 preset 进入实际 child 指标与 zero point 的完整映射为
\begin{center}
\small
\begin{tabular}{c|c|c|c|c}
 类型 & $a_t$ & $a_{0,t}$ & $b_S$ & $b_{S0}$\\ \hline
 h & $a_u+a_v-1$ & $a_{0,u}+a_{0,v}-2\nu_e$ & $b+1$ & $b_0+2\nu_e$\\
 H & $a_u+a_v-1$ & $a_{0,u}+a_{0,v}$ & $b+1$ & $b_0$\\
 massless & $a_u+a_v$ & $a_{0,u}+a_{0,v}$ & $b$ & $b_0$
\end{tabular}
\end{center}
fixed/non-loop line 没有 $b/b_S$ 指标；上表对应的 zero point 改由 sector prefactor 保存。\texttt{shrinkBShift} 与 \texttt{shrinkZeroPointShift} 是由函数系统和 Wronskian 编译出的物理数据，一般不应由用户修改。若确需 override，必须同步重建 child zero point、normalized coefficient 与 parity metadata；额外 normalization 通过 master/basis 选择处理。

\subsection{massless shrink 的缺省吸收与 child-sector parity}

massless theta contact 的 compiled shift 为 $(s_e,z_e)=(0,0)$。对 cycle line，实际映射是
\begin{equation}
 a_t=a_u+a_v,\qquad a_{0,t}=a_{0,u}+a_{0,v},\qquad
 b_S=b,\qquad b_{S0}=b_0.
 \label{v017MasslessDefaultShrink}
\end{equation}
所以没有新增整数幂或 zero-point 幂被吸收。对 fixed/non-loop line，缺省 $B_t-B_s=0$，source/target 保存相同的结构化模长 prefactor，因而 $N_s/N_t=1$；这不表示删除 target 中原有的 fixed-line 因子。normalized contact coefficient 的模长部分为 1，只另乘端点导数的 $-2/+2$ 与共同-theta系数。

massless 的 parity 与 h/H 不同。统一双端点状态的 contact support 为 $n_1+n_2=1$，但式~\eqref{v017MasslessDefaultShrink} 给 $b_S=b$，故
\begin{equation}
 b+n_1+n_2=b_S+1\pmod 2.
 \label{v017MasslessContactParity}
\end{equation}
因此含该 cycle line 的 child-sector affine parity offset 相对 parent 翻转 1。fixed/non-loop massless line 不进入 parity generator，因此不翻转 offset。simultaneous contact 的 offset 按被 shrink 的 massless cycle line 数模 2 累加；该计数与事件只有一个 delta 无关。pure massless family 不自动获得 root parity preset，但 mixed h/H family 若已有认证的 root parity，必须在 massless cycle transition 中保留式~\eqref{v017MasslessContactParity} 的 affine constant。

\subsection{h/H shrink 与 child-sector parity}

裸 H 的 EOM 式~\eqref{bareHEOMindex} 保持 $b+n$ 的模二分级：源 $\{n=2,b\}$ 与三个目标 $\{n=0,b\}$、$\{n=1,b+1\}$、$\{n=0,b+2\}$ 同余；h EOM 也保持同一分级。contact 的分级还取决于实际 shrink 数据。h/H 缺省 Wronskian 分别编译为
\begin{equation}
 (\Delta b,\Delta b_0)_h=(1,2\nu_e),
 \qquad
 (\Delta b,\Delta b_0)_H=(1,0).
 \label{v017HDefaultShrinkParity}
\end{equation}
结合 contact 支持 $n_1+n_2=1$，有
\begin{equation}
 b+n_1+n_2=b+1=b_S\pmod 2,
 \label{v017HContactParity}
\end{equation}
故缺省 h/H shrink 都保持 child parity。h 的 $2\nu_e$ 是 zero point，不是整数 seed 指标；H 的 zero-point shift 为零。fixed/non-loop line 不进入 parity generator。017 必须从实际 compiled \texttt{shrinkTerms}、sector zero-point map 和 contact seed 共同验证式~\eqref{v017HDefaultShrinkParity}，不能只按 preset 文本硬编码。用户修改 child zero point 时，只有相对缺省的变化可精确证明为整数，才按其模二值平移 child offset；否则关闭该 sector 的 parity capability并给出双语诊断。

\end{document}
